% Verbatim copy of /Users/alingalatan/workspaces/grothe/handoff/manuscript/legacy/ch17-crossbound.tex (file dated 2026-10-25). Only this line was added.
% Consolidated source module (formerly Chapter 17). Owner: lower-bound/06-original-responses.tex
\section{Sharper variational certificates and exact-profile constraints (GPT package \code{crossbound-progress}, 15 October 17:22)}
\label{ch:crossbound}
\chapsummary{Variational duals, exact-profile annihilators and coherent tail bounds. Their analytic interfaces remain useful; historical numerical endpoints are not current bounds.}

\paragraph{Provenance (IDE note, 15 October, evening).} This chapter is lower-bound GPT's \code{tex/proposed/ch17-crossbound.tex} from
\pathref{\~{}/Documents/Codex/2026-10-10/plea/outputs/crossbound-progress-20261015/} (108 files, \code{.zip} twin), ported verbatim except for this header and the closing reading note;
the proofs are GPT's and were read, not re-derived line by line, except where the reading note says so. Our replay: \code{MANIFEST.sha256} 108/108; the five proposed edits to
ch1, ch3, ch7, ch15 and \code{main.tex} carried source hashes that matched our files exactly, so they were applied as proposed; \code{verify\_all.py} under \code{python3 -I} from an
isolated directory: all four exact audit suites PASS (1,600 gap/complementarity audits, 2,000 projection-gap identities and 3,000 scalar majorants, the polynomial and exponential
annihilator identities, 102 rational block-PSD certificates); log \pathref{kattack-2026-10-15/gpt\_crossbound\_replay.log}. Status labels are GPT's (\st{P} analytic, \st{C} exact finite
audits, \st{F} quadrature/QMC); the certified bound is unchanged.

\paragraph{Status and purpose.}
The certified master interval reported on 15 October 2026 remains
$1.40160592<\KG<1.404599424\stale{22 Oct}{now $\KG<1.404598161566$}$. The newer unmerged upper branch is reported
at $1.404596071666$ under its analytic hypotheses. [IDE note, 20 Oct: dated; the current Lean interval is $1.401674103\stale{22 Oct}{now $\KG>1.401695322458$}<\KG<1.404599424\stale{22 Oct}{now $\KG<1.404598161566$}$ (\S\ref{sec:certcombinede}); the upper branch \code{work/upper-stable} is not merged as of \code{20eb917}, since it replaces the Lean head computation by a \code{head\_close} hypothesis.] No endpoint is improved
in this chapter. The analytic statements below are \st{P}; the exact finite
audits are supporting \st{C} checks, and all quadrature experiments are
\st{F}. The three additions are a dual for the weaker retained-penalty
problem, zero-projection tests for the actual profiles, and coherent tail
certificates for rounding mixtures.

For the peak cubic row,
$\varepsilon=0.7564$, $L=z(0.8456+0.1432(|z|^2-2))$, and
$S\approx1.37395353$. For a general cubic residual $q$ independent of the
visible $z$, let
$\Lambda_q=\E\Phi_0(|q|)-\E|q|^2/\varepsilon$ and $d_q=S-\Lambda_q$.
The profiles $F_A,F_B$ are the exact charge-one gradients obtained by
averaging the visible coordinate. The target labelled (T$_0$) is
$\norm{\Pc10F_A(q)}^2+\norm{\Pc10F_B(q)}^2
\le0.292(d_q+\eta)$, subject to the stated fibre-majorant scope.

\subsection{Complex primal--dual equality and the exact gap}
\label{sec:complex-duality-proof}
The complex argument was already written out in the 11 October note
\texttt{complex\_duality\_lemma.md}; its audit removes the former
``unwritten'' qualification from the mathematical claim. This is not
a Lean formalization. The standard Gaussian-to-finite-matrix transfer
is proved analytically in the argument below rather than asserted as
a formalized project input. [IDE note: since 15 Oct 17:35 the transfer itself is proved in Lean, \code{lower\_bound\_of\_transferL2}, leaving only the two Gaussian inputs \code{TransferInputsL2}\stale{22 Oct}{proved in Lean: \code{transferInputs\_of\_witnessCap}, \code{witnessCap\_of\_reduction}, \code{lower\_bound\_of\_reduction}}; \autoref{ch:certified}.]

Let $P_k=P_{k+1,k}$ on charge-one complex Gaussian functions and let
\[
N(\lambda)=\sup_n\left\|\sum_{k\ge0}\lambda_kP_k\right\|_{\infty\to1},
\qquad\mathcal Q=\{\lambda:\lambda_0=1,\ |\lambda_k|\le1,
\ \lambda_k\in\R\}.
\]
Let $\mathcal C\subset\ell^1(\R)$ be the convex set of finite mixtures
of coefficient sequences $c_k=\Re\langle P_kf,P_kg\rangle$ for
charge-one disk-valued functions in arbitrary finite dimensions.
Then $\sum_k|c_k|\le1$. Conjugate reflection
$f^\#(z)=\overline{f(\bar z)}$ preserves each charge-one chaos and
conjugates its pairings; averaging a pair with its reflection realizes
the real coefficients as an actual kernel. Disk-valued maps are valid
roundings, or can be decomposed measurably into phase maps.
Circle averaging projects the scalar norm's disk-valued test functions
onto charge one without changing the pairing with $T_\lambda$; hence
$N(\lambda)=\sup_{c\in\mathcal C}\sum_k\lambda_kc_k$.

\begin{theorem}[Complex Gaussian duality, \st{P}]
\label{thm:complex-duality}
\[
\frac1{\KG}=\sup_{c\in\mathcal C}
\left(c_0-\sum_{k\ge1}|c_k|\right)
=\min_{\lambda\in\mathcal Q}N(\lambda).
\]
Finite bounded real Hermite multipliers approximate this minimum.
\end{theorem}
\begin{proof}
The residual rounding transfer gives
$\sup_{c\in\mathcal C}(c_0-\sum_{k\ge1}|c_k|)\le1/\KG$.
Conversely, put $F_n(z)=z/\norm z$ and
$a_n=(\E\norm Z)^2/n\to1$. Its first-chaos vector energy is $a_n$,
and its total vector energy is one. Thus every $\lambda\in\mathcal Q$
has vector payoff at least $2a_n-1$. Approximate $F_n$ by conditional
expectations on finite measurable partitions. The matrix entries are
$\langle1_E,T_\lambda1_F\rangle$, its scalar phase optimum is at most
$N(\lambda)$, and the finite complex Grothendieck inequality gives
$2a_n-1\le\KG N(\lambda)$ on passage to the limit. Only the bounded
$L^2$ operator is required; no pointwise integral kernel is used.
Hence $N(\lambda)\ge1/\KG$.

The real multiplier cube $\mathcal Q$ is compact in the weak-star topology
of $\ell^\infty=(\ell^1)^*$. The pairing with $\mathcal C$, carrying its
$\ell^1$ norm topology, is separately continuous and affine.
Sion's minimax theorem therefore gives
\[
\inf_{\lambda\in\mathcal Q}\sup_{c\in\mathcal C}\sum_k\lambda_kc_k
=\sup_{c\in\mathcal C}\inf_{\lambda\in\mathcal Q}\sum_k\lambda_kc_k
=\sup_{c\in\mathcal C}\left(c_0-\sum_{k\ge1}|c_k|\right).
\]
The preceding two inequalities prove equality. Weak-star lower
semicontinuity gives attainment of the multiplier minimum. There is no
compactness assumption on $\mathcal C$. If $B_d$ restricts the multiplier
to $k\le d$, Ornstein--Uhlenbeck damping and the omitted $L^2$ multiplier
norm give, for $0<r<1$,
\[
\frac1{\KG}\le B_d\le\frac{1}{\KG\sqrt r}+r^{d+1}.
\]
Indeed the damped coefficients are $\lambda_kr^k$, obtained by composing
with the noise operator of parameter $\sqrt r$ and dividing by $\sqrt r$.
Letting $r\uparrow1$ sufficiently slowly as $d\to\infty$ proves convergence.
\end{proof}

For an actual real coefficient mixture define
$D(c)=c_0-\sum_{k\ge1}|c_k|$ and
$W_\lambda(c)=c_0+\sum_{k\ge1}\lambda_kc_k$. The exact identity
\begin{equation}\label{eq:crossbound-gap}
N(\lambda)-D(c)
=[N(\lambda)-W_\lambda(c)]
+\sum_{k\ge1}(|c_k|+\lambda_kc_k)
\end{equation}
splits the gap into nonnegative response and coefficient losses.
Each summand is at least $(1-|\lambda_k|)|c_k|$. Exact complementarity
therefore forces $c_k=0$ at every interior multiplier. In a finite
response-pool LP this statement is exact for the retained head; a valid
upper bound still requires the entire omitted tail. A small pool value
also supplies no global upper bound on the witness norm until all
admissible responses are controlled.

\subsection{Correction: an endpoint Fourier value is not preprocessing-monotone}
\label{sec:damping-correction}
(Same correction as \autoref{ch:problem}, Proposition~\ref{prop:falseceiling}: the diagonal ceiling $1.4045917$ is not a barrier for matrix-covariance preprocessing.)
The inequality $D(c)\le|\sum_k(-1)^kc_k|$ is correct. The additional
claim that preprocessing cannot raise the fixed endpoint value is false.
For a scalar real odd preprocessor, $\rho(i)=ir$ with $|r|\le1$, which
requires the ceiling
\[
C_*:=\sup_{0\le r\le1}\left|\sum_k(-1)^kc_kr^{2k+1}\right|,
\]
rather than its value at $r=1$. For multiple raw coordinates or matrix
preprocessing the supremum must use the corresponding admissible
covariance domain. An endpoint ceiling applies only after proving that
this supremum is attained there.

An exact example is $f(z)=2\sqrt e\,z e^{-2|z|^2}$, with $|f|\le1$.
Its same-map energies are $c_k=4e(k+1)4^k/3^{2k+4}$, whence
\[
C(r)=\frac{4e}{81}\frac{r}{(1+4r^2/9)^2},\qquad
C(1)=\frac{4e}{169}<\frac{e\sqrt3}{72}=C(\sqrt3/2).
\]
The inequality reduces to $3\cdot169^2>288^2$. The example also has a
unimodular-pair realization: write $f=\ph(z)A(|z|)$ and replace $A$ on
the left and right by $A+i\sqrt{1-A^2}S_j$, using separate independent
centered radial signs $S_j$. Their auxiliary chaos terms have zero
cross pairings. Thus the correction is not an artefact of allowing
disk-valued functions. It does not refute the numerical frozen-family
ceiling that explicitly maximizes its diagonal damping variables.

\subsection{A convex dual retaining the first-chaos penalty and cubic consistency}
\label{sec:Kprime-dual}

The diagnostic of subtracting $\norm{Pk_G}^2$ from $\Lambda(G)$ at the
unpenalised phase response is not an upper bound: the response may change
to reduce its penalty. The following dual accounts for this change exactly.
Let $P=P_{1,0}$, $Q=P_{2,1}$, and use
$\langle a,b\rangle=\E[\bar a b]$. For $G\in\operatorname{ran}Q$ and
$v\in\operatorname{ran}P\oplus\operatorname{ran}Q$, set
\[
\mathcal D_G(v):=\E|L+G-v|+\E|L-G+v|+\norm{Pv}^2
+\frac{2\Re\langle Qv,G\rangle-\norm G^2}{\varepsilon}.
\]
Here $\Lambda'(G)$ is the supremum, over unimodular $f,g$ with
$k=(f-g)/2$ and $\varepsilon Qk=G$, of
\[
\Re\E[\overline{L+G}f+\overline{L-G}g]
-\norm G^2/\varepsilon-\norm{Pk}^2.
\]
\begin{proposition}[Retained-penalty dual]
\label{prop:retained-dual}
For every $G$, $\Lambda'(G)\le\inf_v\mathcal D_G(v)\le\Lambda(G)$.
Every individual trial $v$ therefore gives an upper bound, including when
the trials are restricted to a chosen low-dimensional subspace.
\end{proposition}
\begin{proof}
For a feasible pair $f,g$, put $k=(f-g)/2$, so $Qk=G/\varepsilon$.
Subtract its penalised objective from $\mathcal D_G(v)$. The difference is
\[
\begin{split}
&\E\bigl[|L+G-v|-\Re(\overline{L+G-v}f)\bigr]\\
&\quad+\E\bigl[|L-G+v|-\Re(\overline{L-G+v}g)\bigr]
+\norm{Pv-Pk}^2\ge0.
\end{split}
\]
The second inequality follows by taking $v=0$.
\end{proof}

There is also an explicit positive correction. Put
$u_\pm=L\pm G$, $k_0=(\ph(u_+)-\ph(u_-))/2$, and
$r=Pk_0+Qk_0-G/\varepsilon$. On any finite subspace $E$ where the
following integrals are finite, define the quadratic form
\[
\mathcal H_G(v,v)=\norm{Pv}^2+
\frac12\E\left[|v|^2\left(\frac1{|u_+|}+\frac1{|u_-|}\right)\right].
\]
The elementary global inequality
$|u-v|\le |u|-\Re(\overline{\ph(u)}v)+|v|^2/(2|u|)$ gives
\[
\Lambda'(G)\le\Lambda(G)-2\Re\langle v,r\rangle+\mathcal H_G(v,v).
\]
If $H$ is the positive definite Gram matrix of $\mathcal H_G$ in an
orthonormal basis $e_1,\ldots,e_d$ of $E$, and $b_j=\langle e_j,r\rangle$,
then
\[
\boxed{\Lambda'(G)\le\Lambda(G)-b^*H^{-1}b.}
\]
For $E=\operatorname{ran}P$ this becomes
$\Lambda'\le\Lambda-\langle Pk_0,(I+K_{PP})^{-1}Pk_0\rangle$, where
$K$ is the reciprocal-modulus part of $\mathcal H_G$. This first-chaos-only
version does not require cubic consistency; the additional $Q$ directions
do. At a stationary point of $\Lambda$, $Qk_0=G/\varepsilon$, and the
joint correction uses the Schur complement
$I+K_{PP}-K_{PQ}K_{QQ}^{-1}K_{QP}$. It lies between the first-chaos-only
correction and $\norm{Pk_0}^2$.

The exact dual remains valid when reciprocal-modulus integrals diverge;
one must then use $\mathcal D_G$ directly rather than truncate those
integrals without compensation. A rigorous upper matrix enclosure of $H$
gives a conservative correction by inverse order.
If a differentiable minimiser $v$ of $\mathcal D_G$ exists, its phase
response $k_v$ satisfies $Pv=Pk_v$ and $Qk_v=G/\varepsilon$, and the
gap identity above gives equality. Thus the dual is also an exact
constrained response oracle under these stated attainment conditions.

\paragraph{Status.}
This is a dimension-independent analytic reduction, not a proof of the
uniform conjecture $K'$. A global bound for
$\sup_G\inf_v\mathcal D_G(v)$ is still required. The accompanying
composite-radial quadrature probes only the pure junta and are float
diagnostics, not interval certificates or an improved Grothendieck bound.
At a differentiable global maximiser of $\Lambda$, cubic consistency
already holds. Large consistency corrections at a fixed nonstationary
junta therefore do not imply a comparable improvement of the universal
cap; the genuine global improvement here comes from the retained
first-chaos penalty.

\subsection{Exact-profile bounds from functions with zero first-chaos projection}
\label{sec:exact-profile-annihilators}
Let $N$ be the positive Gaussian number operator, with
$Nq=dq$ for a charge-one pure-chaos polynomial of total degree $d\ge3$.
Use the bilinear convention
$\Gamma(f,g)=\sum_j(\partial_j f\,\bar\partial_jg+
\bar\partial_j f\,\partial_jg)$, so
$N(fg)=fNg+gNf-2\Gamma(f,g)$. Put
$U=|q|^2$, $w=\Gamma(q,\bar q)$, $c=\Gamma(q,q)$.
\begin{proposition}[Zero-projection identity, \st{P}]
For a $C^2$ real function $h$ whose composition and derivatives have
polynomial growth, define
\[
\begin{split}
A_h={}&(d-1)qh(U)+2d qUh'(U)\\
&-(\bar q c+2qw)(2h'(U)+Uh''(U))-q^3\bar c h''(U).
\end{split}
\]
Then $A_h=(N-1)[qh(U)]$ and $\Pc10 A_h=0$.
For $d=3$ and $h(U)=e^{-sU}$,
\[
A_s=e^{-sU}\bigl[2q-6s qU+(2s-s^2U)(\bar q c+2qw)
-s^2q^3\bar c\bigr].
\]
\end{proposition}
\begin{proof}
Apply the number-operator chain rule to $F(u)=uh(|u|^2)$:
\[
NF(q)=dqF_u+d\bar qF_{\bar u}-cF_{uu}
-2wF_{u\bar u}-\bar cF_{\bar u\bar u}.
\]
Its derivatives give the displayed expression. Gaussian integration by
parts yields $\E[\bar z_jNf]=\E[\bar z_jf]$ for every coordinate, hence
$\E[\bar z_jA_h]=0$. For the exponential tests all derivatives have
polynomial growth, so Gaussian integrability justifies these operations.
\end{proof}

Choose finitely many rates $s_j>0$ and let $a_0=q$, $a_j=A_{s_j}$.
For arbitrary trial coefficients $\lambda_{Dj}$, projection contraction gives
\begin{equation}\label{eq:annihilator-content}
\sum_{D=A,B}\norm{\Pc10F_D(q)}_2^2
\le\sum_{D=A,B}\E\left|F_D(q)-\sum_j\lambda_{Dj}a_j\right|^2.
\end{equation}
The optimal fixed-input bound is a Schur complement of the finite joint
Gram matrix. Rational trial coefficients give a certificate without solving
the minimisation exactly. The profiles themselves are never differentiated.
This retains the joint law of $(U,w,c)$ and the exact tensor orientation
that the scalar condition (T) discards.

Given the fibre majorant, with $\delta=1/\varepsilon-h_X>0$, a sufficient
condition for the odd-sector target is
\[
\E\left[\Phi_0(|q|)+\frac1\delta\sum_{D=A,B}
\left|F_D(q)-\sum_j\lambda_{Dj}a_j\right|^2\right]
-\frac{\E|q|^2}{\varepsilon}\le S+\eta.
\]
This is a new certificate format, not a universal proof of the condition.
Bounds for the joint weighted integrals, the fibre majorant's scope,
interaction and diagonal uniformity remain obligations.

\paragraph{Diagnostics, \st{F}.}
For rates $1/2,1,2,4,8,16,32,64$, the fitted bound in
\eqref{eq:annihilator-content} is below $0.723$ times
$0.292(d_q+10^{-4})$ at all 27 tested buffered-junta points.
At $\sigma=0.2$, amplitudes $0.3$ and $0.5555$ (normalised fourth moments
$2.3402$ and $5.9993$), the ratios are approximately $0.475$ and $0.722$;
the corresponding optimised scalar-law/cubic-content Gram bounds have
ratios $2.23$ and $1.92$. Higher quadrature confirms these two values,
but the profile interpolation and integral errors are not enclosed.
These measurements therefore prove no new case or endpoint.
The general identities above are analytic; the self-contained exact checker
audits six polynomial cases and the formal exponential chain rule.
On a genuinely coupled two-coordinate cubic, independent grade-three
Gaussian buffering gives normalised fourth moments approximately
$2.2434$, $3.8778$, and $5.8938$. Two independently scrambled QMC probes
give corresponding ratios approximately $0.408$, $0.680$, and $0.708$.
These additional tests also remain \st{F}; sampling and profile errors
are not rigorous enclosures.

\subsection{A coherent tail certificate for response mixtures}
\label{sec:coherent-pool-tail}
The usual tail budget for a mixture sums the separate branch budgets.
It can be sharpened when the omitted tails have certified cross-Gram
matrices. This is an application of the standard trace-norm variational
formula, not a new Grothendieck bound.

Let $P_k$ be orthogonal grade projections, let $(f_i,g_i)$ be a finite
family in a common graded Hilbert space, and let $w_i\geq0$ with
$\sum_iw_i=1$. After retaining grades $k\leq D$, set
\[
 a_i=\sum_{k>D}P_kf_i,\qquad b_i=\sum_{k>D}P_kg_i,
 \quad A_{ij}=\langle a_i,a_j\rangle,\quad
 B_{ij}=\langle b_i,b_j\rangle,\quad W=\operatorname{diag}(w_i).
\]
Then
\begin{equation}\label{eq:coherent-pool-tail}
 \sum_{k>D}\left|\sum_iw_i\langle P_kf_i,P_kg_i\rangle\right|
 \leq\left\|A^{1/2}WB^{1/2}\right\|_*.
\end{equation}
In particular, any Hermitian matrices $X,Y$ satisfying
\[
 \begin{pmatrix}X&W\\W&Y\end{pmatrix}\succeq0
\]
give the square-root-free certificate
\begin{equation}\label{eq:coherent-pool-rational}
 \delta_{\rm tail}\leq\tfrac12\operatorname{tr}(AX+BY).
\end{equation}
Certified Loewner upper bounds for $A,B$ may replace the exact matrices.

For the proof, choose one unit-modulus scalar on each omitted grade to
align its mixed correlation. The resulting grade-diagonal unitary does
not change $A$. Apply the block-positive quadratic form to the transformed
left tail vectors and the right tail vectors. This bounds the sum over
any finite set of grades by the right side of
\eqref{eq:coherent-pool-rational}; monotone passage to all grades and the
trace-norm variational formula give \eqref{eq:coherent-pool-tail}.

The improvement can be strict even for nonsingular Gram matrices. For
two equally weighted branches take
\[
 A=\begin{pmatrix}1&99/101\\99/101&1\end{pmatrix},\quad
 B=\begin{pmatrix}1&-99/101\\-99/101&1\end{pmatrix},\quad W=I/2.
\]
The separate Cauchy--Schwarz budget is $1$. The rational matrices
\[
 X=\frac1{40}\begin{pmatrix}101&-99\\-99&101\end{pmatrix},\qquad
 Y=\frac1{40}\begin{pmatrix}101&99\\99&101\end{pmatrix}
\]
have a positive semidefinite block matrix and certify the sharp budget
$20/101$. All these claims check by exact rational arithmetic.

Optimizing the mixture weights, absolute-value epigraphs for the retained
coefficients, and $X,Y$ is a semidefinite program. Rationalizing its output
leaves a finite PSD certificate and a scalar budget. The additional
analytic inputs are the \emph{joint} tail Gram matrices: storing only the
diagonal energies or separate branch kernels is insufficient. For the
existing scalar preprocessing pool they can be obtained from pairwise
raw-shape correlations in a common graded Gaussian realization, followed
by subtraction of the retained cross-grade Gram matrices.

No improvement of the currently certified constant is asserted here.
The same device may be used for a coherent Taylor or comparison remainder
when its joint Gram data can be enclosed; it cannot recover correlations
already lost in an earlier triangle inequality.

\subsection{What would constitute an improved bound}
The retained-penalty dual is valid for every coupling and diagonal
component, but its supremum over all cubic $G$ has not been bounded by
$S+\eta$. Cubic consistency alone cannot lower a differentiable global
maximum of the unpenalised $\Lambda$, since it already holds at such a
maximum. Large corrections at nonstationary test inputs must therefore
not be advertised as a comparable global gain.

For a buffered Wick input with $\sigma=0.20086016$ and core amplitude
$0.3$, the float raw value is $1.373742781762$, the first-penalty dual
gives $1.373741155969$, and the joint dual gives $1.373705954322$.
Two composite quadratures differ by less than $2\cdot10^{-10}$.
This input was already below $S$, so these numbers certify no new endpoint.

The exact-profile Gram method restores Gaussian constraints absent from
the scalar-law relaxation. Its promising fixed-input fits still require
uniform bounds for weighted $(|q|^2,\Gamma(q,\bar q),\Gamma(q,q))$
integrals over arbitrary cubic tensors. The coherent tail bound requires
joint cross-Gram enclosures for the chosen rounding branches; historical
head tables alone do not contain these data.

The exact supporting suites check 1,600 coefficient-gap identities,
5,000 finite variational/majorant identities, the polynomial and
exponential annihilator formulas, and 102 rational tail certificates.
These checks support the analytic arguments; they do not replace their
arbitrary-dimensional proofs. None changes the formal hypothesis list
of the current certified theorems.

\subsection{IDE reading note}\label{sec:crossbound-reading}
(i) \emph{Zero-projection functions.} That $\Pc10(N-1)F=0$ is immediate ($N$ is multiplication by the chaos degree, so $(N-1)$ kills $\Hc10$); the content of
\S\ref{sec:exact-profile-annihilators} is the explicit closed form of $(N-1)[q\,h(|q|^2)]$ in terms of $(q,w,c)=(q,\Gamma(q,\bar q),\Gamma(q,q))$ of the \emph{same} tensor, which turns
``best approximation of $F_D(q)$ in the orthogonal complement of $\Hc10$'' into a finite Gram problem over explicit functions, without differentiating the profiles. It is an
upper bound on the \emph{actual} content with the tensor orientation kept, i.e.\ a certificate format for (T$_0$) (\autoref{ch:weighted}), and the natural rates $s_j$ are those of the
positive-Laplace representation $F_D(u)=-u\int e^{-s|u|^2}\mu_D(ds)$ of \S\ref{sec:kweighted}. Its reported ratios ($0.47$--$0.72$ of the budget $0.292(d_q+10^{-4})$ on buffered juntas and
a buffered coupled $n=2$ kernel in the kurtosis window, \st{F}) sit where they should: above the exact content (the weighted chain of Theorem~\ref{thm:R} gives $\le0.39\,d_q$ on the
junta and L2 families, \autoref{ch:weighted}) and far below the failing scalar-law chains ($1.9$--$2.2$). What is missing is uniform control of the joint weighted integrals of
$(|q|^2,w,c)$ over all cubic tensors, which is the same obstacle as before in a sharper form.
(ii) \emph{Retained-penalty dual.} We checked the gap identity of Proposition~\ref{prop:retained-dual}: with $v\in\operatorname{ran}P\oplus\operatorname{ran}Q$ and $\varepsilon Qk=G$,
$2\Re\langle v,k\rangle=2\Re\langle Pv,Pk\rangle+2\Re\langle Qv,G\rangle/\varepsilon$, so $\mathcal D_G(v)-J_G(f,g)$ is the two phase gaps plus $\|Pv-Pk\|^2$. Our K$'$ text (\autoref{ch:K} and
\S\ref{sec:kprime-result}) measured $\|Pk_G\|^2$ at the unpenalised best response as a \emph{diagnostic} (16--31\% of the mixed-gain bound; the provably usable part is about $0.41$ of that, \autoref{ch:pinned}, Theorem~\ref{thm:qpenalty}) and already noted that the adversary can
move $f,g$ off the phases; it never used $\Lambda-\|Pk_0\|^2$ as an upper bound, and \S\ref{sec:Kprime-dual} now supplies the rigorous object. The float improvement quoted there
($3.7\cdot10^{-5}$ on a buffered junta already below $S$) is not an envelope gain.
(iii) \emph{Coherent tail certificate.} Standard trace-norm duality, correct. It does not touch the present upper certificate's hypotheses (\autoref{ch:upper}): the Taylor
remainder and the Hardy tail are single-branch items. The place where it could matter is the Cauchy--Schwarz cross term of the comparison step, which is $80\%$ of the certification
allowance $2.73\cdot10^{-6}$ (\autoref{ch:upper}, improvement table); replacing it by signed band contributions is exactly \eqref{eq:coherent-pool-rational}, but that requires certified
\emph{joint} cross-Gram matrices of the omitted tails of the branches, new analytic inputs that the stored head pool does not contain. Lean-relevant only in that sense: a future
component, not a sharpening of an existing \code{KernelHyps}\stale{22 Oct}{\code{KernelHyps} is deleted; the upper bound is \code{fresh\_upper\_rounder} with inputs \code{ComparisonInputs}, \code{energy\_close}, \code{hardy}} field. No numerical gain was measured by GPT or by us.

\section{Why both endpoints stall: a common certificate interface}
\label{sec:forest-reassessment-20261022}
\paragraph{Scope and current numbers (22 October).}
This is a structural reassessment, not a new endpoint. Repository
\code{README.md} at \code{e8d67b5} reports
$1.401691989174\stale{23 Oct}{now $\KG>1.401695322458$}<\KG<1.404599424\stale{22 Oct}{now $\KG<1.404598161566$}$, with the external analytic inputs
documented there. The interval has width $0.002907434826$, about
$0.207\%$ of its upper endpoint. Against the numerical cubic-witness
ceiling $1.402258$ \st{F}, only $19.47\%$ of that width is below the
ceiling; $80.53\%$ is above it. These percentages are diagnostic arithmetic
using a numerical ceiling, not a new rigorous impossibility theorem.
Moving the upper endpoint down to Haagerup's conjectured
$1.404575934663742\ldots$ would remove only $0.81\%$ of this interval.
Thus even completing both present construction-specific targets would
leave a substantial gap. The proved degree-five and degree-seven
obstructions in \autoref{ch:ceilings} supply independent evidence that
the witness degree is a real limitation.

\paragraph{The common object is the attainable response set.}
The exact duality in \S\ref{sec:complex-duality-proof} has no intrinsic
duality gap. Lower-bound relaxations replace the true response set
$\mathcal C$ by an outer approximation; upper constructions provide
an inner collection of actual response pairs. Separate radial laws,
projection energies and mixed contents need not be jointly realizable
by one bounded response. Conversely, a realizable pair with a large
witness payoff need not have a useful rounding residual: its nonlinear
coefficients must cancel in a realizable mixture, including the tail.
This is the precise sense in which a lower-side adversary can help the
upper side. An artificial law passing scalar cuts supplies no such pair.

\paragraph{An exact example of the missing joint information \st{P}.}
Let $X,Y$ be independent standard proper complex Gaussians and
$H(X)=X(|X|^2-2)/\sqrt2$. The pairs $(q,r)=(H(X),X)$ and
$(q,r)=(H(Y),X)$ have identical marginals, unit variances, and the
same zero covariance and pseudocovariance. Nevertheless
\[
 \E[\bar r q|q|^2]=8\sqrt2\quad\hbox{in the first case},
 \qquad \E[\bar r q|q|^2]=0\quad\hbox{in the second}.
\]
Indeed, with $T=|X|^2\sim\operatorname{Exp}(1)$, the first expectation
is $\E[T^2(T-2)^3]/(2\sqrt2)=32/(2\sqrt2)$; the second vanishes
by independence and $\E X=0$. Orthogonality follows from
$\E[T(T-2)]=0$, and pseudocovariance from circular symmetry.
Thus even the complete marginal laws together with their covariance
matrix do not determine the content. This is not a counterexample
to K; it explains why a marginal-law proof needs an additional joint
inequality or a worst-case argument using cubic realizability.

\begin{proposition}[\st{P}: one response pool, two explicit errors]
\label{prop:pool-two-errors-20261022}
Let $c^{(1)},\ldots,c^{(M)}\in\mathcal C$ be actual real coefficient
sequences, and fix a retained level $D$. Write $c(w)=\sum_iw_ic^{(i)}$.
The finite linear-programming value is
\[
 v_D=\min_{\lambda\in[-1,1]^D}\max_i
       \left(c^{(i)}_0+\sum_{k=1}^D\lambda_k c^{(i)}_k\right)
     =\max_{w\in\Delta_M}
       \left(c_0(w)-\sum_{k=1}^D|c_k(w)|\right).
\]
Choose optimizers $\widehat\lambda,w_*$, extend
$\widehat\lambda_0=1$ and $\widehat\lambda_k=0$ for $k>D$.
Suppose the two additional certificates are
\[
 N(\widehat\lambda)\le v_D+\delta,\qquad
 \sum_{k>D}|c_k(w_*)|\le\tau,
 \quad \delta,\tau\ge0,\quad v_D>\tau.
\]
Then
\begin{equation}\label{eq:pool-two-errors-20261022}
 \boxed{\quad \frac1{v_D+\delta}\le\KG\le\frac1{v_D-\tau}.\quad}
\end{equation}
The certified interval for $1/\KG$ has width $\delta+\tau$;
the interval for $\KG$ has width
$(\delta+\tau)/[(v_D-\tau)(v_D+\delta)]$.
\end{proposition}
\begin{proof}
Finite minimax applies to the coefficient cube and the simplex.
Minimizing each $\lambda_k$ gives $-|c_k(w)|$, proving the LP identity.
The actual rounding mixture has full residual at least $v_D-\tau$,
so the rounding transfer gives the upper bound. The globally certified
scalar norm gives the lower bound by the witness transfer. Subtracting
the two reciprocals gives the displayed width. No claim that a
numerically found response maximizes the norm is used.
\end{proof}

\paragraph{Interpretation and a stopping rule.}
The quantity $\delta$ is the certified error of the response maximization
over \emph{all} dimensions and all bounded responses. The quantity
$\tau$ is the tail of the \emph{mixture}, not the sum of unsigned
branch tails unless that is the chosen upper estimate. Solving the
pool LP exactly supplies neither certificate. Without them, $v_D$
has no automatic one-sided relation to $1/\KG$; two opposing
restrictions have been made. Coefficient enclosure and LP rounding
errors must be added explicitly in an implementation.
The natural stopping criterion is a target for $\delta+\tau$, rather
than agreement between two optimizations inside the same pool.

For implementation, exact LP optimality is unnecessary. A feasible
rational mixture with certified head score at least $L$, mixture tail
at most $\tau<L$, and a feasible witness with global norm cap $B$
already gives $1/B\le\KG\le1/(L-\tau)$. This stores checkable
feasible objects rather than trusting an optimizer's termination status.

\paragraph{What complementary slackness actually asks for.}
For $\Delta=N(\lambda)-D(c)$, the earlier gap identity implies
\[
 \sum_{k\ge1}(1-|\lambda_k|)|c_k|\le\Delta.
\]
For a degree-three witness $\lambda_1=-\epsilon$ and
$\lambda_k=0$ for $k\ge2$, this becomes
\[
 (1-\epsilon)|c_1|+\sum_{k\ge2}|c_k|\le\Delta.
\]
Hence a genuinely matching upper mixture must cancel the cubic
coefficient and essentially all higher coefficients. It is the
mixture coefficients that must cancel, not those of every response.
The phase of a linear-plus-cubic polynomial has infinitely many chaos
levels; limiting the witness degree does not remove this upper-side
obligation. This explains why a very good cubic adversary can remain
a poor rounding branch until it is balanced with other responses.

\paragraph{Why complex phases change the real-case proof.}
For real sign responses, $h=(f+g)/2$, $k=(f-g)/2$ satisfy
$h^2+k^2=1$ and $hk=0$. For complex unimodular responses the exact
relations are instead
\[
 |h|^2+|k|^2=1,\qquad \Re(h\bar k)=0.
\]
Both can be nonzero. For disk-valued responses their exact joint
constraint is $|h|^2+|k|^2+2|\Re(h\bar k)|\le1$.
Thus binary exclusion constraints cannot simply be carried over.
For example the CHSH matrix
$M=\left(\begin{smallmatrix}1&1\\1&-1\end{smallmatrix}\right)$
has real scalar norm $2$ and vector norm $2\sqrt2$, but its complex
scalar norm is already $2\sqrt2$: choose the right phases $(1,i)$
and align the left phases with $(1+i,1-i)$. The upper estimate is
$|u+v|+|u-v|\le2\sqrt2$ for unit vectors, by Cauchy--Schwarz and
the parallelogram identity. This Bell witness gives ratio one for
the complex problem. Useful complex orbit constraints must exhibit
a genuine phase-versus-vector gap, which the naive circle cycles lack.

\paragraph{Why neither complex analysis nor a longer alternating pattern is automatic.}
The physical Gaussian kernel is $\sum c_k z|z|^{2k}$, not a holomorphic
function of $z$. Its value at physical correlation $i$ differs from
the analytic continuation $i\sum(-1)^kc_k$ of the real odd series.
A maximum-modulus argument for that continuation does not bound the
rounding objective $c_0-\sum|c_k|$. Moreover,
Corollary~\ref{cor:taper-limit-obstruction} rules out bounded witnesses
converging coefficientwise to the full alternating multiplier: their
norms approach at least one. This does not refute Haagerup's proposed
constant; it refutes that particular route to it. A useful higher-degree
witness must be reoptimized against high-dimensional responses, rather
than obtained by extending a fixed alternating pattern.

\paragraph{Cubics and loss of compactness.}
Charge one permits degrees $1,3,5,\ldots$; the cubic is the first
nonlinear correction and the first mixed three-index tensor. Unlike
one quadratic form, a general cubic cannot be reduced by a unitary
change of basis to independent scalar Wick terms. A shared Gaussian
driver can carry a diffuse quadratic partner. The example
$Q_N=X_1(\sum_{j=1}^N|X_j|^2-N-1)/\sqrt{N+1}$ in the original-response
obstruction ledger retains order-one directional influence while every
fixed-dimensional full cubic projection has vanishing energy.
The limit-law theory of Herry--Malicet--Poly \cite{HMP24} supplies
a structural language for such limits, but not the required quantitative
joint-response comparison at error $10^{-4}$. Marginal convergence of
the cubic alone does not retain its coupling to the same response.

\paragraph{Literature and the practical change of direction.}
HJM's real primal/dual method \cite{HJMreal26} uses actual responses
both to improve the witness and to build its rounding mixture. Its
rigorous certificate uses simpler objects than its numerical search,
so it too separates discovery from a global norm proof. Naor--Regev
\cite{NR14} prove completeness of general real Krivine schemes;
this is not optimality of the original one-dimensional scheme.
The counterexample of BMMN \cite{BMMN13} is a reason to keep
Haagerup's proposed value as a conjecture rather than an optimization
constraint.

The directly useful Bell analogy is membership versus separation:
\href{https://arxiv.org/html/2302.04721v3}{Designolle et al.}
construct local models and separating Bell inequalities with the same
Frank--Wolfe framework. Finding actual response columns is enough to
build a feasible mixture, whereas certifying separation needs the
global scalar optimum. The follow-up \cite{DVP24} makes that
certification bottleneck explicit and notes that its code can treat
the complex problem, though its reported constants concern real-matrix
rank comparisons. Ordinary binary-output XOR games concern the real
Grothendieck constant; their numerical constants cannot be imported here.
Joint moment hierarchies suggest retaining products of the same
response. For our commuting complex variables,
\href{https://arxiv.org/abs/1508.02068}{complex moment/SOS methods}
are a closer finite-dimensional analogue than a noncommutative quantum
hierarchy. Neither automatically supplies an arbitrary-dimensional
Gaussian certificate.

The resulting research priority is to maintain one pool of actual
responses, signed coefficient enclosures and coherent tail Gram data;
reoptimize the witness and rounding mixture against that same pool;
then identify whether $\delta$ or $\tau$ dominates. For the near-term
$1.402$ target, keep the original-response feasibility inequality
\eqref{eq:obstruction-exact-response}. For a substantial closure of
the entire interval, permit higher-degree witnesses and new response
geometry. This is a proposal \st{?}, not a proved rate or endpoint.

\subsection{Implementing the common pool: signed analytic tails and a calibration audit}
\label{sec:shared-pool-implementation-20261022}
\paragraph{Scope (22 October 2026).}
This implements the preceding proposal without claiming a new endpoint.
The repository headline at \code{e8d67b5} is
$1.401691989174\stale{23 Oct}{now $\KG>1.401695322458$}<\KG<1.404599424\stale{22 Oct}{now $\KG<1.404598161566$}$.
The proof, exact replay scripts, actual response parameters and numerical
logs are in
\pathref{handoff/research-third-decimal/shared-response-pool-20261022/},
with companion folders \code{analytic-tail-experiment-20261022},
\code{chirp-tail-experiment-20261022} and
\code{response-pool-audit-20261022}.
The following analytic certificate applies to its stated response family;
it does not settle arbitrary shared cubic responses.

\paragraph{A closed kernel for actual adversaries \st{P}.}
Let $X_j$ be independent proper complex Gaussians assigned odd grades
$w_j$, let $u,v$ be real orthonormal vectors, and define
\[
 F(X)=\operatorname{phase}(u^TX)
       e^{id\Im(\overline{u^TX}\,v^TX)},\qquad
 G(X)=F(s_1X_1,\ldots,s_mX_m),\quad s_j\in\{-1,1\}.
\]
These are bounded charge-one responses in the graded Gaussian model.
Tensor preprocessing realizes their rounding kernels; the corresponding
coefficient sequences also lie in the closure of finite-dimensional
Gaussian response sequences. Thus their coefficients retain joint
realizability rather than merely satisfying moment inequalities.

Put $R(t)=\operatorname{diag}(s_jt^{w_j})$,
$a=u^TRu$, $b=u^TRv$, $k=v^TRv$, and
\begin{align*}
 A&=1+\frac{d^2}{4}(1-a^2-b^2),\\
 B_\pm&=a+\frac{d^2}{4}\{ab^2+k(1-a^2)\}\pm db,\\
 \Delta&=1+\frac{d^2}{2}(1-ak+b^2)
 +\frac{d^4}{16}\{(1-a^2)(1-k^2)-2akb^2+b^4-2b^2\}.
\end{align*}
The full signed correlation series is
\begin{equation}\label{eq:closed-chirp-kernel-20261022}
 K(t)=\frac{\pi B_+}{4\Delta A}
 {}_2F_1\!\left(\frac12,\frac12;2;\frac{B_+B_-}{A^2}\right).
\end{equation}
For a proof, order $(U,U',V,V')$ and tilt its Gaussian covariance
\[
 \Sigma=\begin{pmatrix}1&a&0&b\\a&1&b&0\\0&b&1&k\\b&0&k&1\end{pmatrix}
\]
by the matrix $J$ with entries
$J_{13}=-J_{31}=d/2$ and $J_{24}=-J_{42}=-d/2$.
Then $\det(I-\Sigma J)=\Delta$, and the upper-left block of
$(I-\Sigma J)^{-1}\Sigma$ is
$\Delta^{-1}\left(\begin{smallmatrix}A&B_+\\B_-&A\end{smallmatrix}\right)$.
The two-variable phase integral gives
\eqref{eq:closed-chirp-kernel-20261022}; initially this is an identity
near zero and then an analytic continuation. The determinant and all
four numerators are independently replayed as polynomial identities
over $\mathbb Q[a,b,k,d]$ \st{C}. An independent Wick expansion checks
the first coefficients, including rotations of the full orthogonal
plane of $v$. At $d=0$ the formula reduces to $h(a(t))$.

\begin{proposition}[\st{P}: a signed contour tail]
Suppose $r>1$, $A,\Delta$ have no zeros on $|t|\le r$, and
\[
 \Re\frac{B_+(t)B_-(t)}{A(t)^2}\le1-\delta
 \quad (|t|=r),\qquad 0<\delta\le1.
\]
If the boundary bounds are $|B_+|\le b_{\max}$,
$|A|\ge a_{\min}$ and $|\Delta|\ge\Delta_{\min}$, then
\[
 \sum_{j>N}|[t^{2j+1}]K|
 \le\frac{\pi b_{\max}}{4a_{\min}\Delta_{\min}\sqrt\delta}
       \frac{r^{-2N-3}}{1-r^{-2}}.
\]
\end{proposition}
\begin{proof}
The real-part maximum principle extends the argument bound to the
disk. Euler's integral for ${}_2F_1(1/2,1/2;2;z)$ has modulus at most
$\delta^{-1/2}$ there. Apply Cauchy's estimate and sum only odd degrees.
For rational parameters the contour condition is equivalent to a
rational polynomial inequality in $\cos2\theta$; Bernstein subdivision
certifies the entire circle. Zero-freeness is a separate condition,
proved in the displayed certificates by constant-term domination.
\end{proof}

\paragraph{What cancellation saves \st{P,C}.}
For the actual phase-preprocessor kernel
\[
 P(t)=\frac{1133}{1250}t-\frac{117}{1250}t^3,
 \qquad K(t)=h(P(t)),
\]
the contour certificate at $r=13/10$ gives
$\sum_{j>64}|[t^{2j+1}]K|<1.855\cdot10^{-14}$.
Its unsigned graded energy tail is nevertheless greater than
$0.003838$. Indeed, ordinary total degree is invariant under unitary
rotation, and the weighted degree is at least the ordinary degree.
The unsigned tail therefore dominates
\[
 1-\frac\pi4\sum_{j=0}^{64}a_j>0.003838,
 \qquad a_0=1,\quad
 \frac{a_j}{a_{j-1}}=\frac{(2j-1)^2}{4j(j+1)}.
\]
Both comparisons are exact rational checks using a Machin enclosure
of $\pi$. The factor exceeds $2\cdot10^{11}$. More generally,
conditioning a chirp on $U$ leaves
$\operatorname{phase}(U)e^{-d^2|U|^2/4}$; its normalized Laguerre
coefficients have magnitude asymptotic to $1/(2n)$, so its unsigned
graded energy tail is at least $(1+o(1))/(4N)$.
This explains a specific loss caused by taking absolute values before
combining reflected coefficients. It does not assert that every
response family admits an exponential signed tail. It uses the
holomorphic continuation of the real generating series, not a
holomorphic interpretation of the physical charge-one correlation.

\paragraph{Audit of the numerical pool.}
An active inherited pure-cubic-phase column in the historical
degree-nine search had entries
$-0.7882612315869$ and $-0.0996477635340$, rather than
$-\pi/4$ and $-\pi/32$. The erroneous entries are reproduced by
16-node Gauss--Laguerre quadrature. Polynomial Gram checks do not
detect this phase-response integration error. Repairing recognizable
columns and removing unknown inherited heads changes the floating
restricted LP from $0.712049972007$ to $0.712046050657$ \st{F};
discarding all inherited heads gives the same repaired value.
No certified repository bound used this exploratory pool.

A fresh exchange saves both defining response polynomials, uses the
correct multiplier cube $[-1,1]$, and uses radial Legendre panels.
Its exact-phase calibration errors are below $1.2\cdot10^{-11}$,
but these are not error bounds for arbitrary phase functions.
On refined active contacts, individual coefficients change by up to
$1.65\cdot10^{-4}$ while fixed-witness payoffs change by only a few
$10^{-6}$ \st{F}. Stable payoff is therefore insufficient evidence
for coefficient cancellation. For one saved degree-nine mixture,
the finer index-32 head score is approximately $0.709962$, and the
coherent tail-Gram estimate is approximately $0.009036$, essentially
equal to the triangle estimate. These are numerical diagnostics,
not directed enclosures or a theorem excluding other mixtures.

The analytic chirp pool, enlarged by full orthogonal-plane rotations
and several grades and signs, reaches a floating index-128 head score
$0.7118917430313$. Paying known pure-phase tails gives
$0.7118917280421$, before remaining chirp-tail charges. Its reciprocal
is already about $1.40470799$, weaker than the current upper bound.
This is not a certified ceiling for the continuous response family.
The exact tail replays apply only to their named parameter vectors.
A further floating search with general Hermitian quadratic phase
modulation, five exchange rounds and 594 retained columns reaches
$0.7118923203268$ after known pure-phase tail charges, before the
remaining tails. This also fails to improve the current upper bound;
it does not exclude asymmetric response pairs or a better search.

\paragraph{The resulting certificate discipline.}
For a nominal mixture head $\widehat c$, certified head error
$\epsilon$ in $\ell^1$, and complete tail $\tau$, the upper denominator
is
\[
 \widehat c_0-\sum_{j=1}^N|\widehat c_j|-\epsilon-\tau.
\]
The lower endpoint still needs a separate global estimate
$N(\lambda)\le B$. Realizability, coefficient accuracy, complete tails
and the global response oracle are different obligations.
The present implementation supplies a closed analytic route through
the first three for a restricted family. It neither proves
Conjecture K nor supplies the missing global norm cap. The next
search must optimize the complete signed coefficient objective,
rather than select responses solely for their low-degree payoff.

% Verbatim copy of /Users/alingalatan/workspaces/grothe/handoff/manuscript/legacy/global-response/spherical-dual-transfer.tex (file dated 2026-10-25). Only this line was added.
\subsection{An exact spherical dual with a uniformly priced operator tail}
\label{sec:spherical-dual-transfer-20261022}

\paragraph{Scope and status (22 October).}
This is a global analytic reduction \st{P}, with exact supporting checks
\st{C} and scalar tail/cap assembly \st{L}. It preserves the same bounded
responses in every Gaussian dimension. It is not a new numerical endpoint
or a claim of literature novelty. A useful global spherical norm cap
remains to be proved; the small truncation allowance below is not the
error of the entire spherical relaxation.

Write $P_k=P_{k+1,k}$ and
$N_G(\lambda)=\sup_n\sup_{|f|,|g|\le1}
\sum_k\lambda_k\Re\langle P_kf,P_kg\rangle$ on charge-one responses.
On the probability sphere $S^{2m-1}\subset\mathbb C^m$, let $\Pi_r$
project onto the spherical harmonics of bidegree $(r+1,r)$, and let
$N_{S,m}(a)$ be the analogous norm of $\sum_r a_r\Pi_r$.

\begin{theorem}[\st{P}: global response-level spherical transfer]
\label{thm:spherical-dual-transfer}
Suppose $\lambda_k=0$ for $k>D$ and $m\ge2D+3$. Define
\[
 C_{m,r}=\binom{m+r-2}{r},\qquad
 q_{m,r}(s)=\frac{P_r^{(m-2,1)}(2s-1)}{C_{m,r}},\qquad
 a_{m,r}=\sum_{j=0}^D\lambda_j[s^j]q_{m,r}(s).
\]
Then $a$ is bounded and
\begin{equation}\label{eq:spherical-dual-global}
 \boxed{N_G(\lambda)\le N_{S,m}(a).}
\end{equation}
If $m>2D+3$, the spherical operator is compact and its norm is attained.
Explicitly,
\begin{equation}\label{eq:spherical-dual-inverse}
 a_{m,r}=\frac1{C_{m,r}}\sum_{j=0}^{\min(D,r)}
 \lambda_j(-1)^{r-j}\frac{(r+m)_j}{j!}\binom{r+1}{j+1}
 =O(r^{2D+3-m}).
\end{equation}
\end{theorem}
\begin{proof}
The normalized zonal harmonics are $\Phi_{m,r}(z)=zq_{m,r}(|z|^2)$.
Their dimensions are
\[
 d_{m,r}=\binom{m+r}{r+1}\binom{m+r-1}{r}
          -\binom{m+r-1}{r}\binom{m+r-2}{r-1},
\]
where the second term is zero for $r=0$. Orthogonality gives
\begin{equation}\label{eq:spherical-dual-b}
 z|z|^{2k}=\sum_{r=0}^k b_{m;k,r}\Phi_{m,r}(z),\qquad
 b_{m;k,r}=\frac{d_{m,r}(m-1)!k!(k+1)!}{(k-r)!(m+k+r)!}.
\end{equation}
Indeed the coefficient is $d_{m,r}(m-1)$ times the integral of
$s^{k+1}(1-s)^{m-2}q_{m,r}(s)$ on $[0,1]$; shifted Jacobi Rodrigues
and $r$ integrations by parts evaluate that integral as
$\binom{k}{r}B(k+2,m+r-1)/C_{m,r}$. Thus $b\ge0$ and
$\sum_r b_{m;k,r}=1$.

For independent Gaussian vectors $X_1,\ldots,X_m\in\mathbb C^n$, set
\[
 F_X(u)=f\Bigl(\sum_j u_jX_j\Bigr),\qquad
 G_X(u)=g\Bigl(\sum_j u_jX_j\Bigr).
\]
They remain disk-valued and charge one, and use the \emph{same} matrix
$X$. Mehler and the trace on the spherical harmonic space yield
\[
 \E_X\Re\langle\Pi_rF_X,\Pi_rG_X\rangle
 =\sum_{k\ge r}b_{m;k,r}c_k,\qquad
 c_k=\Re\langle P_kf,P_kg\rangle.
\]
Apply $\ell(p)=\sum_{j=0}^D\lambda_j[s^j]p$ to
\eqref{eq:spherical-dual-b}: for \emph{every} $k$,
$\sum_{r\le k}b_{m;k,r}a_{m,r}=\lambda_k$.
The endpoint Jacobi derivative formula gives
\eqref{eq:spherical-dual-inverse}. Boundedness of $a$, stochastic
row sums and $\sum_k|c_k|\le1$ justify rearrangement, proving
\eqref{eq:spherical-dual-global}. When $a_r\to0$, finite-rank
truncations converge in $L^2$ operator norm. Weak-star compactness of
the charge-one unit disk and compactness of the operator give attainment.
This makes no coordinate-core or independent-driver assertion.
\end{proof}

\begin{corollary}[\st{P,C}; scalar truncation logic \st{L}]
\label{cor:spherical-dual-cubic-tail}
For $\lambda=(1,-\epsilon,0,\ldots)$, $m\ge6$, $0\le\epsilon\le1$, put
\[
 a_{m,r}=(-1)^ru_{m,r},\qquad
 u_{m,r}=\frac{r+1}{\binom{m+r-2}{r}}
       \left(1+\frac{\epsilon r(r+m)}2\right).
\]
Let $a^{[L]}$ be $a$ through index $L$ and zero afterward. Then
\begin{equation}\label{eq:spherical-dual-cubic-cap}
 \boxed{N_G(1,-\epsilon)\le N_{S,m}(a^{[L]})+u_{m,L+1}.}
\end{equation}
No omitted response-energy-tail hypothesis is required.
\end{corollary}
\begin{proof}
For $r\ge1$ the ratio is
\[
 \frac{u_{m,r+1}}{u_{m,r}}=
 \frac{r+2}{m+r-1}
 \frac{2+\epsilon(r+1)(r+m+1)}{2+\epsilon r(r+m)}\le1.
\]
The cross-multiplied difference factors as
\[
 2(m-6)+6(1-\epsilon)+
 \epsilon\{(m-5)(r-1)(r+m+2)+(m-6)(m+1)\}\ge0.
\]
Bessel gives $\sum_r|\langle\Pi_rF,\Pi_rG\rangle|\le1$;
hence the truncation error is at most $u_{m,L+1}$ uniformly.
The truncation concerns the \emph{dual operator}, so the slow-phase-tail
obstruction in Corollary~\ref{cor:joint-cubic-tail-obstruction} still holds.
\end{proof}

\paragraph{The outer relaxation has a separate, explicit error.}
For sphere responses $F,G$, their actual Gaussian radial lifts
$f(Z)=F(Z/\norm Z)$, $g(Z)=G(Z/\norm Z)$ have coefficients
\[
 c_k=\sum_{r\le k}v_{m;k,r}d_r,\quad
 d_r=\Re\langle\Pi_rF,\Pi_rG\rangle,\quad
 v_{m;r+j,r}=\gamma_{m,r}
 \frac{(r+1/2)_j^2}{j!(m+2r+1)_j},
\]
\[
 \gamma_{m,r}=\frac{\Gamma(m+r+1/2)^2}
 {\Gamma(m)\Gamma(m+2r+1)}.
\]
Expand in the radial Laguerre basis
$R^{2r+1}L_j^{(m+2r)}(R^2)$ to prove this formula; Gauss's
hypergeometric identity gives column sums one. It preserves mixed
signed coefficients. For the cubic game the lifted multipliers are
$\mu_0=\gamma_{m,0}[1-\epsilon/(4(m+1))]$,
$\mu_1=-\epsilon\gamma_{m,1}$ and $\mu_r=0$ for $r\ge2$.
Comparison with $a$ gives
\begin{equation}\label{eq:spherical-dual-outer-error}
 0\le N_{S,m}(a)-N_G(1,-\epsilon)\le
 \max\left\{\frac{1+\epsilon}{4m},\,
 \frac{2(1+\epsilon)}{m-1}+\frac{9\epsilon}{4m},\,
 \frac{6[1+\epsilon(m+2)]}{m(m-1)}\right\}.
\end{equation}
Here $1-\gamma_{m,r}\le(2r+1)^2/(4m)$: Wallis and
$\gamma_{m+1,0}/\gamma_{m,0}=1+1/[4m(m+1)]$ give the initial
deficit $1/(4m)$; the ratio
$(m+r+1/2)^2/[(m+2r+1)(m+2r+2)]$ has deficit at most
$2(r+1)/m$. Sum the deficits. Full Laguerre calculations and the
comparison proof are in the accompanying proof file.

\paragraph{Calibration and computational obstruction.}
For $F=G=\ph(u_1)$ the exact rational spherical energies are
\[
 e_r=d_{m,r}I_r^2,\quad
 I_0=\frac{2^{2m-1}m!(m-1)!}{(2m)!},\quad
 I_{r+1}=-I_r\frac{2r+1}{2m+2r+1}.
\]
They satisfy $\sum_ra_{m,r}e_r=\pi/4-\epsilon\pi/32$,
since every nonzero restriction of the Gaussian phase is a unitary
rotation of $\ph(u_1)$. This is a normalization test, not optimality.

At $\epsilon=149/200$, $m=32,L=12$ the truncation allowance is
$17513/209010560960<8.38\cdot10^{-8}$. However, the complex rank is
\[
 R_{m,L}=\binom{m+L}{L+1}\binom{m+L-1}{L},\qquad
 R_{32,12}=796315556846635874048.
\]
The truncated norm has a maximizing pair of phases of polynomials in
this finite range, but direct coefficient enumeration is not credible.
Furthermore, \eqref{eq:spherical-dual-outer-error}, not this tiny tail,
prices the finite-$m$ outer relaxation. A useful continuation must
exploit unitary symmetry or a positive operator certificate. No such
norm certificate, improved endpoint or exhaustive response
classification is asserted here.

There is also an intrinsic dimension floor. If $a_0=1$, then
\[
 N_{S,m}(a)\ge\frac1{K_G^{\mathbb C}(m)}.
\]
The vector response $u\mapsto u$ has unit norm and exact payoff one,
since each coordinate is a first spherical harmonic. Approximate it
by unit-valued step functions on small sphere cells and apply the
finite-dimensional Grothendieck inequality to the resulting matrix.
Its scalar optimum is at most $N_{S,m}(a)$; circle averaging deals
with non-equivariant step phases. Hence a target below this floor is
impossible for that $m$, however accurately the operator tail is
computed. This explains why a small-$m$ tail example alone supplies
no route to the current endpoint.

\paragraph{An exact obstruction to a squared-kernel certificate.}
There is a stronger warning than the finite coefficient count. The
dimension formula and the inverse multipliers give the exact identity
\[
 \boxed{d_{m,r}|a_{m,r}|^2=
 \frac{(m+2r)(m+r-1)(r+1)}{m-1}
 \left(1+\frac{\epsilon r(r+m)}2\right)^2.}
\]
Indeed $d_{m,r}/C_{m,r}^2=(m+2r)(m+r-1)/[(m-1)(r+1)]$.
For $\epsilon>0$ the right side grows as
$\epsilon^2r^7/[2(m-1)]$, and for $\epsilon=0$ it grows cubically.
Thus the compact operator is not Hilbert--Schmidt in any of these
dimensions and has no square-integrable product-sphere kernel.
Its truncated squared Hilbert--Schmidt norm grows as
$\epsilon^2L^8/[16(m-1)]$, even while its omitted operator norm
decays as $O(L^{5-m})$. A squared-kernel estimate therefore cannot
supply the required signed norm certificate for this operator.
This does not invalidate Hilbert--Schmidt estimates for different
comparison operators elsewhere in the manuscript.
More generally, if $\lambda_D\ne0$ is the highest coefficient,
\[
 d_{m,r}|a_{m,r}|^2\sim
 \frac{2\lambda_D^2}{(m-1)(D!)^2((D+1)!)^2}\,r^{4D+3}.
\]
Increasing dimension improves the individual multiplier decay, but
does not cure this multiplicity obstruction. Increasing witness degree
makes the divergent exponent larger.

\paragraph{Sources, replay and formalization.}
The harmonic input is classical; see
\href{https://arxiv.org/abs/2411.12579}{Defant et al.},
\href{https://dlmf.nist.gov/18.37}{DLMF 18.37} and
\href{https://arxiv.org/abs/1903.07697}{Peterson--Sengupta}.
The objective and completeness motivation are those of
\cite{HJMreal26,NR14}. The project-specific inverse and error ledger
are derived above. Full proofs, exact verifier, droplet log and the
source map of all 125 current manuscript TeX files are in
\pathref{/Users/alingalatan/workspaces/grothe/outputs/spherical-dual-transfer-20261022/}.
The module
\pathref{/Users/alingalatan/workspaces/grothe/Grothe/LowerAnalytic/SphericalTransfer.lean}
proves the all-index scalar tail and signed cap-assembly logic.
The analytic bridge and a spherical norm cap remain outside that
formalization; none is added to the master endpoint's hypotheses.

% Verbatim copy of /Users/alingalatan/workspaces/grothe/handoff/manuscript/legacy/global-response/binding-audit-20261023.tex (file dated 2026-10-25). Only this line was added.
\subsection{Reconstructing the binding certificate configurations}
\label{sec:binding-audit-20261023}
\paragraph{Scope (23 October).}
This audit uses master \code{98f61ec} and its lifted joint-moment
certificate at $\varepsilon=149/200$. It changes neither endpoint.
The finite data are reconstructed exactly; the four radial countermodels
below concern a fixed table-and-cut relaxation, not actual responses.
The report, data and logs are in
\pathref{/Users/alingalatan/workspaces/grothe/outputs/binding-audit-20261023/};
the account of the discarded variables is in
\pathref{/Users/alingalatan/workspaces/grothe/docs/BINDING\_AUDIT\_2026-10-23.md}.

\paragraph{The finite envelope \st{C}.}
The read-only prefix of the existing data generator reproduces all
21 batches, 351 nodes and 8,756 boxes. Their strings agree exactly with
the committed Lean data; all box bounds, node ranges and the partition
are replayed. This does not recompute the source Gaussian integrals or
independently recheck every half-line node majorant.
The maximum box bound is between $0.713421798525$ and $0.713421798526$,
leaving about $1.4745\cdot10^{-9}$ below the headline $0.7134218$.
The explicit pad is $10^{-11}$. Neither is the principal loss.
There are 130 boxes within $10^{-7}$ of that maximum, with enclosing
$t$ interval $[0.84425,0.8545]$. Adaptive refinement to a common target
does not identify these with near-maximizing actual responses.

For rule 7234 the certificate expression is maximized near
$(t,\tau,x)=(0.8495,0.245542324,0.7265625)$ and has the decomposition
\[
 -0.719702191279+0.055407255822+1.408882874811
       -0.031166140839+10^{-11},
\]
with the displayed terms rounded. They are respectively the visible
quadratic, beta completion, moment-majorant expectation, pinned price
and pad. They must not be treated as numerical error terms.
Selecting direct chords from the available nodes of the same potentials
saves about $2.9500\cdot10^{-6}$ at this point. Across twelve selected
rules the pointwise savings range from zero to about
$1.3985\cdot10^{-5}$. No replacement cover is asserted.

\paragraph{Exact feasible laws for a fixed relaxation \st{C}.}
Let $\widehat\Psi_i$ denote the stored source interpolants, and fix the
profile, price and parameter point of a selected rule. Define
\[
 R_{\rm tab}(\nu)=t^2-2\bar\ell t+\frac{2\bar\beta^2}{\varepsilon}
 +\int\sum_iw_i\widehat\Psi_i(\delta_i\tau\sqrt u)\,d\nu(u)
 -D\tau^2+10^{-11}.
\]
Four rational atomic probability laws satisfy exactly
\[
 \mathbb EU=1,\quad \mathbb EU^2=2+x^2,\quad
 \mathbb EU^3\le6+9x^2+40.45313x^3,\quad
 \mathbb E\sqrt U\le0.973612157,
\]
and all 671 listed cut functions on their certified ranges. This bank
contains the 306 spectral/joint cuts, 126 distinct inherited mixed Stein
cuts, and 239 additional distinct source Stein cuts from the five cells
referenced by the four rules. It is not the set of all possible tests.
The exact strict lower bounds for $R_{\rm tab}$ are
\[
\begin{array}{c|c|c}
 \text{rule}&\text{atoms}&R_{\rm tab}(\nu)>\\\hline
 7234&10&0.713412812507\\
 5661&15&0.713414640742\\
 1247&12&0.713406413945\\
 8701&11&0.713352277645
\end{array}
\]
and each exceeds $500/701$. Floating LPs only discover the laws.
The independent standard-library verifier recomputes rational moments,
every cut and the table score without an optimizer.

Thus a dual majorant using this fixed table and finite cut bank cannot
have cost below the corresponding feasible primal score. The generic
finite separation argument is kernel-proved in
\code{Grothe.Lower.PrimalAudit}; these four numerical instances remain
external exact checks. No claim of literature novelty is made.
In particular the laws are not Gaussian cubic laws or bounded response
pairs. The tables are upper envelopes of the exact potentials.
These witnesses therefore obstruct this specified relaxation only;
they neither violate the original inequality nor exclude improved
tables, stronger cuts, reoptimized row profiles, or a different witness
such as $1891/2500$. No minimax floor over all profiles is proved.

\paragraph{Which joint data the next reduction must retain.}
The pre-supremum reduction has the same response pair, remote cubic $p$,
gate vector $(p,X_2,X_3)$ and $Qh$. After mixing, it drops
\[
 \varepsilon\|Qh+(\bar\beta/\varepsilon)v\|^2,
 \qquad
 \mathbb E\|Y+C^{-1}h\,p\|_C^2,\quad Y=(X_2,X_3).
\]
The beta mixture saving has already been used. At rule 7234, vanishing
of the latter remainder would require approximately
$X_2=-0.532438p$, $X_3=-0.256700p$.
The saved radial data specify these target ratios but not the actual
remainders or the response-weighted fields $P(fA_i)$ and $P(gB_j)$.
Those fields must be retained and constrained before the fibre
supremum. A free cross block can be set to zero, so appending block
positivity after forgetting it supplies no new restriction.

There is also a row-consensus loss. If $F_i(f,g)$ is the fibre
expression before its supremum, then
\[
 \sup_{f,g}\sum_iw_iF_i(f,g)
 \le\sum_iw_i\sup_{f,g}F_i(f,g).
\]
The left side averages the quadratic matrices, visible terms and
coupling vectors before maximizing over one common section. The
current scalar tables use the right side. The nonnegative difference
has not been quantitatively certified here. It is a further joint
observable, not a proved numerical improvement or a second allowance
to spend an existing square.

% Verbatim copy of /Users/alingalatan/workspaces/grothe/handoff/manuscript/legacy/global-response/joint-response-lift-20261023.tex (file dated 2026-10-25). Only this line was added.
\subsection{A shared Gram lift before the fibre supremum (23 October)}
\label{sec:shared-response-gram-lift}

\paragraph{Scope.}
The following extraction is dimension-free and applies to every actual
disc-valued response pair. The stability constraint additionally requires a
globally near-maximizing pair, as in
\S\ref{sec:fresh-second-chaos-stability}. Neither the extraction nor the
floating-point diagnostics below gives a new numerical endpoint. In particular,
the atomic laws of \S\ref{sec:binding-audit-20261023} do not specify the new
observables and cannot be declared excluded by assigning them convenient values.

\paragraph{The common projection budget.}
Use the realification of complex Gaussian $L^2$, with scalar product
$\operatorname{Re}\mathbb E\overline u v$. Let $P$ be the full first-chaos
projection, and let $P_0$ project onto any chosen complex subspace of that
chaos. Fix bounded real features $A_1,\ldots,A_r,B_1,\ldots,B_s$ and form the
\emph{single} tuple
\[
 v=(fA_1,\ldots,fA_r,gB_1,\ldots,gB_s),\qquad
 x_i=P_0v_i,\quad z_i=v_i-x_i.
\]
Write $\operatorname{Gr}(u)_{ij}=\operatorname{Re}\langle u_i,u_j\rangle$.
The measured matrix $X=\operatorname{Gr}(x)$ and the residual matrices
\[
 V=\operatorname{Gr}(z)=\operatorname{Gr}(v)-X,\qquad
 K=\operatorname{Gr}(Pv-x)
\]
satisfy
\begin{equation}\label{eq:shared-gram-interval}
 \boxed{0\preceq K\preceq V,\qquad
       \operatorname{Gr}(Pv)=X+K.}
\end{equation}
Indeed $Pz_i=Pv_i-x_i$. For any real $c$,
$c^{\mathsf T}Kc=\|P\sum c_i z_i\|^2$ lies between zero and
$\|\sum c_i z_i\|^2=c^{\mathsf T}Vc$. Orthogonality to the measured
space proves the other two identities. No independence, conditional
Gaussianity, tensor-rank bound or estimate of an individual Hermite tail is used.

Crucially, the mixed block of $\operatorname{Gr}(v)$ is
\[
 \operatorname{Gr}(v)_{i,r+j}
   =\mathbb E\bigl[\operatorname{Re}(\overline f g)A_iB_j\bigr].
\]
It is a moment of the \emph{same pair}. Separate Bessel budgets discard it.
The dimension of \eqref{eq:shared-gram-interval} is $r+s$, regardless of the
number of unmeasured Gaussian directions.

\paragraph{A sharp scalar completion screen.}
For one feature on each branch, put $c=X_{12}$,
$v=V_{12}$ and $R=\sqrt{V_{11}V_{22}}$. Every possible residual coefficient
$k=K_{12}$ obeys the sharper interval
\begin{equation}\label{eq:shared-cross-interval}
 \boxed{|2k-v|\le R.}
\end{equation}
One proof uses the reflection $2P-I$: it is an isometry and
$2\langle Pz_1,Pz_2\rangle-\langle z_1,z_2\rangle
=\langle(2P-I)z_1,z_2\rangle$. Apply Cauchy--Schwarz.
For $V=I_2$, projection onto $(1,1)$ attains $k=1/2$; projection onto
$(1,-1)$ attains $k=-1/2$. Thus the radius $R/2$ cannot be reduced from
Gram information alone. The separate estimate $|k|\le R$ has twice that radius
when $v=0$.

Suppose stability gives $|c+k|\le a$. Eliminating \emph{all} completions yields
\begin{equation}\label{eq:shared-completion-screen}
 \boxed{|2c+v|\le R+2a.}
\end{equation}
This follows by subtracting $2k-v$ from $2(c+k)$ and using the triangle
inequality. For example, $c=1/5$, $v=0$, $R=1/10$ and $a=13/100$ fail
the screen: $2/5>9/25$. Separate Cauchy and stability still allow $k=-2/25$.
This is an exact example separating the two completion tests, not an example
at a Grothendieck binding node.

\paragraph{Connection to actual response stability.}
Retain the weights and measured moments before taking any fibre supremum:
\[
 D_f[i,j]=\mathbb E(|Tg|A_iA_j),\qquad
 D_g[i,j]=\mathbb E(|Tf|B_iB_j),\qquad
 C=(X+K)_{fg}.
\]
If the pointwise feature budgets are at most $L^2$, a globally
$\delta$-near-maximizing pair has the necessary finite-dimensional completion
\begin{equation}\label{eq:shared-stability-completion}
 0\preceq K\preceq V,\qquad
 \begin{pmatrix}D_f&\varepsilon(X+K)_{fg}\\
 \varepsilon(X+K)_{gf}&D_g\end{pmatrix}
       +10L^2\sqrt\delta\,I\succeq0.
\end{equation}
For scalar features, one may use
$a=\sqrt{(D_f+10L^2\sqrt\delta)(D_g+10L^2\sqrt\delta)}/\varepsilon$
in \eqref{eq:shared-completion-screen}. The signed weights
$\operatorname{Re}(\overline f Tg)$ and
$\operatorname{Re}(\overline g Tf)$ give the stronger original Hessian
constraint. They must not silently be substituted for the positive weights
in a computation that assumes the latter.

\paragraph{Finite dual prices for the residual block.}
For real symmetric matrices $Z,Y$, the two finite checks
$Y\succeq0$ and $Y-Z\succeq0$ imply
\begin{equation}\label{eq:shared-gram-price}
 \langle Z,K\rangle\le\langle Y,V\rangle
 \qquad(0\preceq K\preceq V).
\end{equation}
Here $\langle A,B\rangle=\operatorname{tr}(AB)$. Indeed
$\langle Z,K\rangle\le\langle Y,K\rangle\le\langle Y,V\rangle$.
Rational Gram factorizations of $Y$ and $Y-Z$ suffice to check the prices;
an optimizer's eigenvalues are unnecessary. The rank-one version is included
in the Lean module below.

This budget is exact for arbitrary Hilbert-space feature vectors: given
$0\preceq K\preceq V$, choose $A=V^{\dagger/2}KV^{\dagger/2}$ on
$\operatorname{ran}V$. Then $0\preceq A\preceq I$ and the block operator
\[
 \Pi=\begin{pmatrix}A&\sqrt{A(I-A)}\\
                      \sqrt{A(I-A)}&I-A\end{pmatrix}
\]
is an orthogonal projection. Applying it to the columns of
$(V^{1/2},0)^{\mathsf T}$ realizes $K$. This does \emph{not} realize
arbitrary such columns as $fA_i,gB_j$ on Gaussian space. The remaining
structure must enter through their local response moments.

\paragraph{A moment-constrained fibre support interface.}
There is a dimension-free way to choose the features so that their response
moments can be replaced by moments of the two polynomial images, with a
controlled near-maximizer error. This avoids declaring arbitrary bounded
responses independent of those images.

\paragraph{An image-determined near-maximizer lift.}
Let $a=|Tg|$, $b=|Tf|$, $\phi_g=\operatorname{phase}(Tg)$ and
$\phi_f=\operatorname{phase}(Tf)$. A globally $\delta$-near-maximizing pair
satisfies
\begin{equation}\label{eq:shared-weighted-alignment}
 \mathbb E\,a|f-\phi_g|^2\le2\delta,\qquad
 \mathbb E\,b|g-\phi_f|^2\le2\delta.
\end{equation}
For example, the first bound follows from the pointwise disc inequality
$a|f-\phi_g|^2\le2(a-\operatorname{Re}\overline f Tg)$ and
$\mathbb E|Tg|\le M$, since replacing $f$ by $\phi_g$ is admissible in
the same dimension. Both inequalities hold for the original pair; successive
phase replacements or existence of a maximizer are unnecessary.

Choose real bounded image-dependent rules $\alpha_i,\beta_j$ with
$\sum\alpha_i^2,\sum\beta_j^2\le1$, fix $\tau>0$, and set
\[
 A_i=\sqrt{\frac a{a+\tau}}\,\alpha_i,
 \qquad B_j=\sqrt{\frac b{b+\tau}}\,\beta_j.
\]
Then, summing over features,
\begin{equation}\label{eq:shared-cutoff-alignment}
 \sum_i\|(f-\phi_g)A_i\|_2^2\le\frac{2\delta}{\tau},\qquad
 \sum_j\|(g-\phi_f)B_j\|_2^2\le\frac{2\delta}{\tau}.
\end{equation}
Indeed $\sum A_i^2/a\le1/\tau$, including $a=0$ by continuity.
Define the image tuple
$\widetilde v=(\phi_g A,\phi_f B)$. Its full, measured and residual
Gram matrices differ in operator norm from those of $v=(fA,gB)$ by at most
\[
 4\sqrt{2\delta/\tau}.
\]
To see this, the synthesis operators of both tuples have Hilbert--Schmidt
norm at most $\sqrt2$, and their difference has norm at most
$2\sqrt{\delta/\tau}$. Use
$\|U^*U-\widetilde U^*\widetilde U\|
\le(\|U\|+\|\widetilde U\|)\|U-\widetilde U\|$.
The same argument applies after any orthogonal projection.
For the cross block alone the sharper error is
$2\sqrt{2\delta/\tau}$, because each individual branch has synthesis norm
at most one. Consequently the \emph{image-determined} necessary condition is
\begin{equation}\label{eq:shared-image-stability}
 \boxed{\begin{pmatrix}D_f&\varepsilon\widetilde C\\
       \varepsilon\widetilde C^{\mathsf T}&D_g\end{pmatrix}
       \succeq-\left(10+2\varepsilon\sqrt{2/\tau}\right)\sqrt\delta\,I,}
\end{equation}
where $\widetilde C=\operatorname{Gr}(P\widetilde v)_{fg}$ and $D_f,D_g$
are unchanged. Its measured and unmeasured matrices satisfy the exact common
Gram interval for $\widetilde v$. For fixed positive $\tau$, the error vanishes
uniformly in the dimension. This is a necessary condition on global
near-maximizers, not a bound on all arbitrary cubic perturbations.

At fixed polynomial-image state the features and $\widetilde v$ are now
known functions, rather than moments of free response variables. One possible
rule uses clipped $\operatorname{Re}(\overline{\phi_g}\ell_i)/(a+\eta)$,
normalized to the declared feature budget, and its other-branch counterpart.
The image law must still come from the same pair under $T$; arbitrary pairs
of polynomial images are an outer relaxation. No claim is made that the
present radial tables already enforce that compatibility.

\paragraph{The local moment support.}
Here is a way to retain those moments rather than append free matrices after
the old table. Fix a remote state $w$, including the conditional polynomial
images $(Tf)(\cdot,w),(Tg)(\cdot,w)$, and any bounded feature rule. For
fixed real quadratic-dual variable $y$, let $F_y(z;u,v)$ be the linear
two-response payoff from the old fibre identity and let
$m(z;u,v)$ collect the measured coefficients, the joint Gram entries and the
weighted feature moments just defined. If their fibre expectation is $b$, then
for \emph{every} real multiplier vector $\theta$,
\begin{align}
 \sup_{\substack{|f|,|g|\le1\\\mathbb E_z m(z;f,g)=b}}
 \{\mathbb E_zF_y(z;f,g)-y^{\mathsf T}My\}
 &\le\mathbb E_z\sigma_z(y,\theta)-y^{\mathsf T}My-\theta\cdot b,
 \label{eq:shared-local-support}\\
 \sigma_z(y,\theta)
 &=\sup_{|u|,|v|\le1}
       \{F_y(z;u,v)+\theta\cdot m(z;u,v)\}.\nonumber
\end{align}
This is just the pointwise support inequality, integrated before subtracting
the measured moments. Supremizing in $y$ preserves the bound; infimizing
over chosen prices remains valid. No tensor supremum is interchanged with an
infimum. Each inner support is on two complex discs, hence four real variables.
It retains $\operatorname{Re}(\overline u v)$; replacing it by two independent
one-disc supports would discard precisely the joint information.

For response-independent features the moment integrands give the quadratic
support
\[
 c+\sup_{|u|,|v|\le1}
 \{a|u|^2+b|v|^2+\operatorname{Re}(\overline p u+\overline q v)
                         +2\operatorname{Re}(\overline u wv)\}.
\]
One disc can be eliminated exactly. For $r\ge0$ define
\[
 \chi_b(r)=\begin{cases}b+r,&2b+r\ge0,\\
                         -r^2/(4b),&2b+r<0.\end{cases}
\]
Maximizing first the phase and then the modulus of $v$ gives the equivalent
\emph{two-real-variable} support
\begin{equation}\label{eq:shared-two-real-support}
 c+\sup_{|u|\le1}\{a|u|^2+\operatorname{Re}(\overline p u)
                   +\chi_b(|q+2\overline w u|)\}.
\end{equation}
Indeed $\chi_b(r)=\max_{0\le t\le1}(bt^2+rt)$; its maximizer is $1$ in
the first branch and $-r/(2b)$ in the second. Both the bound and attainment
are kernel-checked in the Lean module. Negative diagonal prices must use the
second branch when applicable: assuming optimal responses are unimodular
would otherwise give a false support value. If $a,b\ge0$, separate convexity
puts both responses on their circles and reduces the last search to one angle.
These simplifications do not apply unchanged to features depending nonlinearly
on the response, such as clipped response-weighted scores.

The standalone checker \code{tools/joint\_disk\_support.py} implements a
rational enclosure of \eqref{eq:shared-two-real-support}. On a rectangle $R$
meeting the unit disc, let $n_-$ be its minimum squared distance to the origin
and $n_+=\min\{1,\max_{\rm corners}|u|^2\}$. Bound the quadratic term by
$an_+$ if $a\ge0$, and by $an_-$ otherwise. The linear term is bounded by
the minimum of $|p|$ and its maximum over the four corners. Finally use
\[
 |q+2\overline w u|\le
 \min\{|q|+2|w|,\max_{\rm corners}|q+2\overline w u|\}.
\]
The corner bound follows from convexity of the norm; $\chi_b$ is increasing
because it is the supremum of $bt^2+rt$ with $t\ge0$. Adding these three
bounds gives a valid cell upper bound. Feasible rational points give lower
bounds. Rational square-root enclosures follow from integer square roots;
dyadic subdivision preserves a cover of the disc, and cells are discarded
only when disjoint from it or below the incumbent lower bound. Thus the
returned interval is rigorous even if a cell budget is exhausted before
reaching the requested tolerance. This checker covers fixed rational local
coefficients; it does not integrate over Gaussian states or certify a norm cap.

\paragraph{An exact $3\times3$ complex SDP for the local support.}
There is a stronger simplification of the same quadratic support. Put
\[
 H=\begin{pmatrix}
 0&\overline p/2&\overline q/2\\
 p/2&a&w\\q/2&\overline w&b
 \end{pmatrix}.
\]
Then, exactly,
\begin{align}
 \sup_{|u|,|v|\le1}(1,u,v)^*H(1,u,v)
 &=\max_{X\succeq0,\ X_{00}=1,\ X_{11},X_{22}\le1}
           \operatorname{tr}(HX)\label{eq:shared-exact-local-sdp}\\
 &=\min_{\substack{\lambda_1,\lambda_2\ge0\\
             \operatorname{diag}(\lambda_0,\lambda_1,\lambda_2)-H\succeq0}}
            (\lambda_0+\lambda_1+\lambda_2).\nonumber
\end{align}
Here $X$ is complex Hermitian and $\lambda_0$ is real without a sign
restriction. This is a local two-response result, not an exact SDP for the
whole Grothendieck problem.

Equivalently, the complete local quadratic moment body is
\[
 \boxed{\operatorname{conv}\{zz^*:z=(1,u,v),\ |u|,|v|\le1\}
       =\{X\succeq0:X_{00}=1,\ X_{11},X_{22}\le1\}.}
\]
Thus a loss in this part of a lifted certificate can be removed without a
higher moment hierarchy. This is local convexification only; it does not
identify which fields of such moments arise from a common Gaussian response
pair with prescribed polynomial projections.

For completeness, the feasible set of $X$ is compact and convex. If an extreme
point has rank $r\ge2$, write $X=RR^*$ with $R$ of full column rank. The real
space of Hermitian $r\times r$ matrices has dimension $r^2\ge4$. Therefore
there is a nonzero Hermitian $U$ with
$\operatorname{diag}(RUR^*)=0$, since this imposes at most three real equations.
For sufficiently small $t>0$, the distinct matrices
$R(I\pm tU)R^*$ are positive semidefinite and preserve all three diagonals.
Their average is $X$, contradicting extremality. Every extreme point thus has
rank one. Since $X_{00}=1$, write it as $(1,u,v)(1,u,v)^*$ after a common phase
rotation, with $|u|,|v|\le1$. A linear objective attains its optimum at an
extreme point, proving the first equality. The strictly feasible point
$\operatorname{diag}(1,1/2,1/2)$ and a sufficiently large diagonal dual matrix
give strong duality and attainment for the displayed primal and dual.

This is the usual complex correlation-matrix rank argument, here adapted to
the two-disc support; see Li--Tam, \emph{A note on extreme correlation matrices},
SIAM J. Matrix Anal. Appl. 15 (1994), 903--908,
\href{https://doi.org/10.1137/S0895479892240683}{doi:10.1137/S0895479892240683}.
The real symmetric perturbation space has dimension $r(r+1)/2$ instead;
the argument permits rank two with three real diagonal constraints. Thus this
exact local simplification specifically uses complex response phases. Adding
more response branches also invalidates the rank-one dimension count.

For rational data, any rational feasible dual prices are accepted by checking
the seven principal minors of the $3\times3$ Hermitian slack, together with
$\lambda_1,\lambda_2\ge0$. A floating SDP only proposes the prices; a small
rational diagonal loading makes the slack certifiably positive when necessary.
\code{tools/two\_disc\_sdp.py} contains a standard-library exact verifier and
a separate optional proposal mode. Its four local examples, including an
interior coupled optimizer, have exact dual--primal gaps below $10^{-8}$.
This replaces the rectangle mesh for the quadratic local support. The rectangle
checker remains an independent enclosure method. Neither checker supplies the
Gaussian integral or the global joint-image law cover.

There is also a direct continuous-cell rule. If
$H(s)=\sum_j\omega_j(s)H_j$, with nonnegative weights summing to one, interpolate
feasible dual prices by the same weights. The slack is the corresponding
convex combination of the certified positive semidefinite slacks, so the local
support is at most $\sum_j\omega_j(s)\sum_i\lambda_{j,i}$ throughout the cell.
This applies directly when the coefficient matrix is affine in the cell
parameters; more general inputs need a certified coefficient enclosure.
The old negative quadratic in $y$ remains separate and is spent once.

For a fixed lifted state the dual support is convex in the moment prices:
it is a supremum of affine functions of those prices, even after supremizing
in $y$. This does not make the inner maximization in $y$ concave or justify
exchanging the tensor supremum and price infimum.

For a positive definite $M$, the outer quadratic dual is coercive. Bounded
features and fixed polynomial images give an integrable bound
$C_0+C_1\|y\|-\lambda_{\min}(M)\|y\|^2$ on its value, with an added
explicit constant for the moment prices. Thus a finite ball in $y$ and the two
discs suffice for interval subdivision; Gaussian polynomial tails are paid as
in the old potential checker. This gives a finite support obligation for each
specified lifted state and price. It does \emph{not} yet give a finite cover
of all admissible remote laws of those states. In particular the conditional
images must remain the actual images of the same responses, rather than free
polynomials chosen separately for each moment.

\paragraph{A bilinear alternative for fixed feature rules.}
There is also a useful way to avoid the absolute-weight price when the real
features $A,B$ are response-independent. Let $M_A$ denote multiplication by
$A$ and use the signed Hessian form from
\S\ref{sec:fresh-second-chaos-stability}. For $\lambda\ge0$,
\[
 J(f,g)+\lambda H(A,B)
 =\operatorname{Re}\langle f,S_\lambda g\rangle,\qquad
 S_\lambda=T+\lambda(M_{A^2}T+TM_{B^2}+2\varepsilon M_APM_B).
\]
If the features have a uniform bound in every dimension, the near-maximizer
Hessian theorem and a maximizing sequence imply
\[
 \boxed{\|T\|_{\infty\to1}\le
        \sup_n\|S_\lambda\|_{\infty\to1}.}
\]
Indeed, at a $\delta$-near-maximizer the right payoff is at least
$\|T\|-\delta-10\lambda L^2\sqrt\delta$; now let $\delta\downarrow0$.
The same argument permits finite nonnegative sums of bounded feature tests.
This is a new certificate interface, not a claim that the modified norm is
smaller or already bounded. Response-dependent clipped scores do not define
one fixed bilinear operator and must instead use the lifted support interface.

\paragraph{A common-row fallback.}
For row functionals $F_i(f,g)$ before their fibre supremum,
$\sup_{f,g}\sum w_iF_i\le\sum w_i\sup_{f,g}F_i$.
Since the pre-supremum functional is affine in $(M,L,\delta a)$, its common
potential uses $\overline M=\sum w_iM_i$, $\overline L=\sum w_iL_i$ and
$\overline A=\sum w_i\delta_i a_i$; its dual subtracts
$m^{\mathsf T}\overline M^{-1}m$. Averaging $M_i^{-1}$ would be incorrect.
This replacement preserves the existing common-square mixing identity.
It cannot also claim the old full discarded-square deficit a second time.

\paragraph{Quadratic features close the missing first-chaos map exactly.}
The free completion in \eqref{eq:shared-stability-completion} is unnecessary
for an important feature class. For a Hermitian matrix $M$, let
\[
 A_M(z)=z^*Mz-\operatorname{tr}M,\qquad
 W_{jl;i}={:}z_jz_l\overline z_i{:}.
\]
The pointwise Wick identity is
\begin{equation}\label{eq:shared-quadratic-wick}
 A_Mz_l=\sum_jM_{lj}z_j+\sum_{i,j}M_{ij}W_{jl;i}.
\end{equation}
Consequently every disc-valued response, not only a near-maximizer, satisfies
\begin{align}
 [P(fA_M)]_l
 &=\sum_j\overline{M_{lj}}\langle z_j,f\rangle
       +\sum_{i,j}\overline{M_{ij}}\langle W_{jl;i},f\rangle
       \nonumber\\
 &=\sum_j\overline{M_{lj}}\langle z_j,Tf\rangle
       -\varepsilon^{-1}\sum_{i,j}\overline{M_{ij}}
                                      \langle W_{jl;i},Tf\rangle.
       \label{eq:shared-quadratic-closure}
\end{align}
Here $[Ph]_l=\langle z_l,h\rangle$. Reality of $A_M$ moves it to the test
probe, and \eqref{eq:shared-quadratic-wick} lies entirely in the first and
$(2,1)$ chaoses. Thus higher response coefficients vanish from the pairing.
No multiplication operator on all of $L^2$ is assumed: the bounded response
times $A_M$ is in $L^2$, and the probe product is a fixed polynomial.

Write the last expression as $\mathcal D_M(Tf)_l$. The full cross block is
now exactly
\[
 C_{ij}=\operatorname{Re}\langle\mathcal D_{M_i}(Tf),
                                      \mathcal D_{N_j}(Tg)\rangle.
\]
It is a contraction of the actual linear and cubic tensors. Changing its
unmeasured part independently while leaving those images fixed would violate
\eqref{eq:shared-quadratic-closure}. This addresses precisely the completion
freedom exposed by the score-feature diagnostic; it does not assert that
every quadratic-feature stability matrix is violated at a binding state.

The features are unbounded, so the bounded-feature remainder constant cannot
be reused. There is an explicit dimension-free substitute. If
$\|A_M\|_2=\|M\|_{\rm HS}=a$, then
\[
 \mathbb EA_M^4\le9a^4,\qquad
 \mathbb EA_M^6\le265a^6,\qquad
 \mathbb EA_M^8\le14833a^8.
\]
To prove these, diagonalize $M$ and write
$A_M=\sum\lambda_j(|z_j|^2-1)$. For $k\ge2$ its cumulant is
$(k-1)!\sum\lambda_j^k$, whose absolute value is at most $(k-1)!a^k$.
The moment--cumulant partition formula has nonnegative coefficients; replacing
each cumulant by this bound gives the corresponding even moment of
$a(\operatorname{Exp}(1)-1)$. Those centered even moments are the derangement
numbers $9,265,14833$.

For two such features with norms $a,b$, the $L^8$ remainder in
\S\ref{sec:fresh-second-chaos-stability}, using $\|Tf\|_2,\|Tg\|_2\le1$,
is at most
\[
 C(A,B)\le\frac{183}{4}(a^4+b^4)+\frac{45}{2}a^2b^2
                   +\frac{489}{20}(a^3b+ab^3)
             \le48(a^2+b^2)^2.
\]
Here $\sqrt{14833}<122$, $\sqrt{265}<163/10$ and $\varepsilon\le1$.
For the last inequality put $S=a^2+b^2$, $r=ab$; the difference multiplied
by $276$ is
$(138r-489S/20)^2+(9279/400)S^2\ge0$.
Therefore orthonormal Hermitian feature families give the necessary matrix
inequality
\begin{equation}\label{eq:shared-quadratic-stability}
 \boxed{\begin{pmatrix}D_f&\varepsilon C\\
                         \varepsilon C^{\mathsf T}&D_g\end{pmatrix}
             \succeq-16\sqrt{3\delta}\,I}
\end{equation}
for global $\delta$-near-maximizers. This also holds with the original signed
weights. No clipping error or free residual Gram completion is present. The
pointwise Wick identity, Hilbert-space probe closure and rational remainder
budget are kernel-checked in \code{QuadraticFeatures.lean}; the Gaussian
moment--cumulant application and matrix assembly here are written proofs.

The constant feature may be included without losing closure. For
$A=c+A_M$ with $\|A\|_2=a$, Minkowski and Cauchy--Schwarz give
$\|A\|_4\le2a$, $\|A\|_6\le\sqrt8\,a$ and
$\|A\|_8\le\sqrt{13}\,a$. Indeed $9=(\sqrt3)^4$,
$265<343=(\sqrt7)^6$ and $14833<20736=(\sqrt{12})^8$.
Since $(\sqrt8)^3<23$, the remainder becomes at most
\[
 \frac{507}{8}(a^4+b^4)+40a^2b^2+
                         \frac{69}{2}(a^3b+ab^3)
                 \le67(a^2+b^2)^2.
\]
For the last inequality, multiplying the difference by $347$ gives
$(347r/2-69S/2)^2+(541/8)S^2\ge0$.
Thus affine-quadratic feature families have matrix error at most
$4\sqrt{67\delta}\,I$. This second rational budget is also in Lean.

Thus the next uniform certificate can retain joint polynomial images and
their explicit quadratic contractions, use the exact local two-disc SDP for
moment prices, and bound the weighted image integrals. It still needs a
global joint-tensor inequality and complete row cover. A marginal radial law
alone contains neither the contraction map nor those weighted integrals.

\paragraph{Formalization and diagnostic status.}
The Hilbert-space Gram interval, its nested-projection specialization,
reflection estimate, scalar completion screen, factored prices, finite
moment-support assembly, pointwise alignment and cutoff inequalities,
cross-inner-product perturbation, scalar disc support and its attainment,
and near-maximizer price accounting are formalized in
\code{Grothe/LowerAnalytic/ResponseGram.lean}. The Gaussian application,
Gaussian image-lift and matrix-norm error estimates, the two-disc support
construction, matrix dilation and bilinear operator
identity above are paper proofs; they are not asserted to have been formalized
by those abstract lemmas. Discovery scripts and independent-grid results are
recorded in \code{docs/JOINT\_RESPONSE\_LIFT\_2026-10-23.md}.
No numerical endpoint is changed.

\code{TwoDisc.lean} additionally proves the positive-slack factorization and
its universal support bound, and sharp bounds and attainment for the four local
examples. These proofs use rational sum-of-squares factorizations, so the four
sharp local support values, unlike the floating SDP proposals, require no
external numerical hypothesis. The general rank-one SDP exactness argument
and the principal-minor proposal checker remain paper and external checks.

% Verbatim copy of /Users/alingalatan/workspaces/grothe/handoff/manuscript/legacy/global-response/common-stein-kernel-20261023.tex (file dated 2026-10-25). Only this line was added.
\subsection{Common Stein kernels and a genuine-cubic floor for the frozen tables}
\label{sec:common-stein-reassessment-20261023}

\paragraph{Scope and reassessment.}
The finite-bank countermodels of \S\ref{sec:binding-audit-20261023} do not
establish a barrier for the complete radial Stein relaxation. In fact all four
can be excluded by merging two inherited tests before completing their squares.
Reoptimizing the law largely replaces those atoms without recovering a useful
payoff gain. A stronger obstruction is supplied below by an explicit genuine
finite-dimensional cubic and all nineteen frozen source tables. Neither result
changes a numerical Grothendieck endpoint.

\paragraph{One kernel for every test.}
For a unit genuine $(2,1)$ cubic retain the existing definitions
\[
 U=|q|^2,\qquad X=\Gamma(q,\bar q)/3-1,\qquad
 r=\Re(\Gamma(q,q)\bar q^2)/(3U),
\]
with $r=0$ at zero. Write $P=\E UX$, $Q_c=\E Ur$ and
$\chi=\E U^2-2$. The inherited inputs are
\[
 \E X=0,\quad \E X^2\le2P/3,\quad \E r^2\le2Q_c/3,
 \quad 2P+Q_c=\chi,\quad Q_c\le\chi/3.
\]
For each continuous piecewise-affine test $\phi_j$, eventually zero, set
$\Phi_j(u)=\int_0^u\phi_j(s)\,ds$ and
$d_j(u)=\phi_j(u)-\Phi_j(u)/u$, with $d_j(0)=0$. All tests use the same
$X,r$ in
\[
 \E[\Phi_j(U)-\phi_j(U)]
   =\E[X\phi_j(U)+rd_j(U)].
\]
These are existing Gaussian inputs, proved in
\code{Grothe/LowerAnalytic/RadialCut.lean}; no conditional Gaussianity is added.

\begin{proposition}[Exact finite-feature compression, paper proof]
For a fixed law with $\E U^2<\infty$, put
\[
 v_X=(1,U,\phi_1(U),\ldots,\phi_s(U)),\qquad
 v_r=(U,d_1(U),\ldots,d_s(U)),
\]
$G_X=\E v_Xv_X^{\mathsf T}$, $G_r=\E v_rv_r^{\mathsf T}$, and
$A_j=\E[\Phi_j(U)-\phi_j(U)]$. Existence of a common square-integrable
kernel for these finite Stein equations and separate budgets is equivalent to
existence of $P,Q_c,\xi,s_X,s_r$ satisfying
\[
 \begin{pmatrix}G_X&(0,P,\xi)^{\mathsf T}\\(0,P,\xi)&s_X\end{pmatrix}\succeq0,
 \qquad
 \begin{pmatrix}G_r&(Q_c,A-\xi)^{\mathsf T}\\(Q_c,A-\xi)&s_r\end{pmatrix}\succeq0,
\]
\[
 s_X\le2P/3,\quad s_r\le2Q_c/3,\quad
 2P+Q_c=\chi,\quad Q_c\le\chi/3.
\]
\end{proposition}
\begin{proof}
Necessity follows from the two augmented Gram matrices. Conversely, a positive
block $\left(\begin{smallmatrix}G&y\\y^{\mathsf T}&s\end{smallmatrix}\right)$
implies $(c^{\mathsf T}y)^2\le s c^{\mathsf T}Gc$. The functional
$\sum c_i v_i\mapsto c^{\mathsf T}y$ is therefore well-defined and bounded
on the feature span. Riesz representation supplies one vector in that span
with all the prescribed pairings and squared norm at most $s$. Apply this to
each block. The same $\xi$ makes every Stein equation hold simultaneously.
The argument includes singular Grams; positivity enforces their nullspace
compatibility. Orthogonal projection preserves all retained moments and lowers
both variance costs, so nonnegative-price coupled variance budgets can also be
retained.
\end{proof}

This finite-family equivalence does not encode the extra pointwise gradient
cone $X\ge-1$, $|r|\le1+X$, characterize cubic laws, or impose bounded-response
realizability. Its matrices are determined by the law and the chosen tests;
they are not arbitrary completions of unknown Gaussian projection data.
The simultaneous projection principle is formalized as
\code{Grothe.Lower.SteinFusion.common\_kernel\_projection}.

\paragraph{An exact compiler for old tests.}
For an inherited cut with $\mu=\lambda/2$, $t=-\lambda/6$, $\nu=0$,
divide all parameters and test values by $\lambda$. It becomes
\[
 G_i=f_i^2+2g_i^2+\Phi_i-\phi_i+(u^2-2)/3,
 \qquad f_i=-(\phi_i+\beta_i)/2,\quad g_i=-d_i/2.
\]
Merge the two tests and centering parameters with weight $0<w<1$, refining
their meshes to the exact union. Then
\begin{equation}\label{eq:common-stein-fusion}
 G_w=wG_1+(1-w)G_2
 -w(1-w)\bigl((f_1-f_2)^2+2(g_1-g_2)^2\bigr).
\end{equation}
The merged test is admissible under the existing
\code{Grothe.Lower.spec\_of\_inputs} theorem. Thus its expectation is
nonnegative, even when separate nonnegativity of the old cuts would allow
incompatible predictions for their common kernel. More generally,
\[
 \frac{x^2}{a}+\frac{y^2}{b}-\frac{(x+y)^2}{a+b}
 =\frac{(bx-ay)^2}{ab(a+b)},\qquad a,b>0.
\]
Both identities, the associated inequalities and the specialization of
\eqref{eq:common-stein-fusion} to the actual spectral-cell formula are
Lean proved. This is an exact convex compilation of existing inputs, not a
new universal chaos inequality.

\paragraph{Exact and exploratory checks.}
Each of the four old atomic laws passes all 671 old cuts. A merge of two of
those inherited cuts has expectation below $-0.0661$, $-0.2629$, $-0.0989$ and
$-0.0122$ at indices $7234,5661,1247,8701$, respectively; both constituent
expectations remain nonnegative. The stored rational certificates verify
the union mesh, mixing weight, negative expectation and variance-saving
identity. Four independently proposed 16-knot tests also separate the laws.
All eight exact replays and eight corruption controls pass.

This does not quantify a comparable objective improvement. Three adaptive
tests change the four finite-grid scores by approximately $0$, $2.36\cdot10^{-8}$,
$1.17\cdot10^{-9}$ and $7.51\cdot10^{-9}$: other atomic laws replace them.
A joint law/kernel conic discovery solve at rule 5661 gives approximately
$0.713412131532$, against $0.713414640743$ for the old grid. This is a
3,735-point calculation, not a continuous-law cap. The compressed SDP attempts
failed numerically; no certificate is inferred from them.

\begin{proposition}[Genuine-cubic floor for all frozen row mixtures, \st{P,C}]
Let $f_i$ be the nineteen frozen source upper tables of master
\code{98f61ec}, retaining their exact parameters $\ell_i,\beta_i,\delta_i,
k_i,h_i,C_i$. Let $w_i\ge0$, $\sum_iw_i=1$, and average these parameters
before forming the beta square and Schur price. At
\[
 \epsilon=149/200,\qquad t=341/400,\qquad
 \tau=239818283/10^9,
\]
there is an explicit unit genuine cubic $q_{128}$ on 8,384 remote coordinates
such that every such mixture satisfies
\begin{equation}\label{eq:genuine-frozen-table-floor}
 \boxed{t^2-2t\ell_w+\frac{2\beta_w^2}{\epsilon}
 +\sum_iw_i\E f_i(\delta_i\tau|q_{128}|)
 -(k_w-h_w^{\mathsf T}C_w^{-1}h_w)\tau^2
 >\frac{14267}{20000}=0.71335>\frac{500}{701}.}
\end{equation}
Adding the existing positive padding only increases the left side.
\end{proposition}
\begin{proof}
For independent proper complex coordinates $X_i,Y_{ij}$, $1\le i\le j\le m$,
put
\[
 q_m=\frac{\sum_iX_i^2\bar Y_{ii}
       +\sqrt2\sum_{i<j}X_iX_j\bar Y_{ij}}{\sqrt{m(m+1)}}.
\]
The two coordinate blocks are disjoint, so this is a genuine Wick cubic.
Conditionally on $X$, its variance is $R_m^2/[m(m+1)]$,
where $R_m=\sum|X_i|^2\sim\Gamma(m,1)$. Hence
\[
 q_m\stackrel{\mathrm{law}}=\frac{R_mG}{\sqrt{m(m+1)}},\qquad
 \E|q_m|^{2k}=\frac{k!(m)_{2k}}{[m(m+1)]^k},
\]
with $G$ independent and standard proper complex Gaussian. Its variance is one
and its kurtosis excess is $(8m+12)/(m(m+1))$. Take $m=128$; the exact mean
modulus and sixth-moment gates are checked from these formulas.

For each table, choose slopes $0\le d\le s_{\max,i}$ and intercepts
$b=\min_j(y_j-dx_j)$. Each line $b+ds$ lies below every table cell and the
entire final ray. Their maximum is a convex minorant. If its initial value
and slope are $b_0,d_0$ and its nonnegative slope jumps at $a_j\ge0$ are
$\Delta_j$, its exact Rayleigh expectation is
\[
 H(\sigma)=b_0+\frac{\sqrt\pi\sigma}{2}
       \left(d_0+\sum_j\Delta_j\operatorname{erfc}(a_j/\sigma)\right).
\]
As a function of $R_m$, this expectation is convex. Partition its gamma
radius into $[0,64]$, width-two intervals through 192, and $[192,\infty)$.
Conditional Jensen bounds the expectation below by the sum of cell masses
times $H$ at the exact conditional mean. Both tails are retained. The masses
and first moments follow from
\[
 \Pr(R_m>a)=e^{-a}\sum_{j=0}^{m-1}\frac{a^j}{j!},\qquad
 \E[R_m\mathbf1_{R_m>a}]=m\Pr(R_{m+1}>a).
\]
The verifier recomputes these quantities and the error functions using
256-bit Arb balls.

For fixed rational $\beta_0,z$, completing squares gives
\[
 \frac{2\beta_w^2}{\epsilon}\ge
 \frac{4\beta_0\beta_w-2\beta_0^2}{\epsilon},\qquad
 h_w^{\mathsf T}C_w^{-1}h_w\ge
 2z^{\mathsf T}h_w-z^{\mathsf T}C_wz.
\]
All nineteen $C_i$ are checked positive definite over the rationals. The
resulting lower bound is affine in $w$; checking nineteen directed strict
inequalities therefore proves it for every mixture. The table JSON is
hash-pinned and agrees exactly with the previously audited source rows.
Discovery weights or quadrature convergence are not proof premises.
\end{proof}

\paragraph{The obstruction and the next target.}
This is a floor for the \emph{stored upper-table expression}, not for the
exact fibre potentials or the original witness norm. We have not shown that
the same response pair produces $q_{128}$ at these $t,\tau$. It is not a
counterexample to the original-response conjecture. It does not apply to new
profiles, shifted potentials, an averaged potential formed before its supremum,
or another value of $\epsilon$.

It does establish that arbitrary improvements to radial constraints alone,
even all valid constraints of genuine cubic laws, cannot finish this frozen
nineteen-table architecture: the explicit cubic remains admissible. The next
global certificate must change the profiles or retain the common response
relation $p=Q_w\E_Zk$ together with the actual images $Tf,Tg$ and the negative
projection costs. Rejecting individual atomic laws is no longer an adequate
measure of progress.

\paragraph{Files and formal scope.}
The report is \code{docs/COMMON\_STEIN\_REASSESSMENT\_2026-10-23.md}; proof
data and logs are under \code{outputs/common-stein-20261023/}.
\code{tools/verify\_common\_stein.py} uses only the standard library;
\code{tools/verify\_genuine\_table\_floor.py} additionally uses
\code{python-flint}. The six general fusion and simultaneous projection
results are Lean formalized in \code{Grothe/Lower/SteinFusion.lean}.
The finite-feature Schur equivalence and the Gaussian/Jensen application are
paper proofs with external exact/directed numerical replays. No headline
endpoint is changed.

